% Self-contained ICML-style named-author white paper.
\begin{filecontents*}[overwrite]{algorithm.sty}
% ALGORITHM STYLE -- Released 8 April 1996
%    for LaTeX-2e
% Copyright -- 1994 Peter Williams
% E-mail Peter.Williams@dsto.defence.gov.au
\NeedsTeXFormat{LaTeX2e}
\ProvidesPackage{algorithm}
\typeout{Document Style `algorithm' - floating environment}

\RequirePackage{float}
\RequirePackage{ifthen}
\newcommand{\ALG@within}{nothing}
\newboolean{ALG@within}
\setboolean{ALG@within}{false}
\newcommand{\ALG@floatstyle}{ruled}
\newcommand{\ALG@name}{Algorithm}
\newcommand{\listalgorithmname}{List of \ALG@name s}

% Declare Options
% first appearance
\DeclareOption{plain}{
  \renewcommand{\ALG@floatstyle}{plain}
}
\DeclareOption{ruled}{
  \renewcommand{\ALG@floatstyle}{ruled}
}
\DeclareOption{boxed}{
  \renewcommand{\ALG@floatstyle}{boxed}
}
% then numbering convention
\DeclareOption{part}{
  \renewcommand{\ALG@within}{part}
  \setboolean{ALG@within}{true}
}
\DeclareOption{chapter}{
  \renewcommand{\ALG@within}{chapter}
  \setboolean{ALG@within}{true}
}
\DeclareOption{section}{
  \renewcommand{\ALG@within}{section}
  \setboolean{ALG@within}{true}
}
\DeclareOption{subsection}{
  \renewcommand{\ALG@within}{subsection}
  \setboolean{ALG@within}{true}
}
\DeclareOption{subsubsection}{
  \renewcommand{\ALG@within}{subsubsection}
  \setboolean{ALG@within}{true}
}
\DeclareOption{nothing}{
  \renewcommand{\ALG@within}{nothing}
  \setboolean{ALG@within}{true}
}
\DeclareOption*{\edef\ALG@name{\CurrentOption}}

% ALGORITHM
%
\ProcessOptions
\floatstyle{\ALG@floatstyle}
\ifthenelse{\boolean{ALG@within}}{
  \ifthenelse{\equal{\ALG@within}{part}}
     {\newfloat{algorithm}{htbp}{loa}[part]}{}
  \ifthenelse{\equal{\ALG@within}{chapter}}
     {\newfloat{algorithm}{htbp}{loa}[chapter]}{}
  \ifthenelse{\equal{\ALG@within}{section}}
     {\newfloat{algorithm}{htbp}{loa}[section]}{}
  \ifthenelse{\equal{\ALG@within}{subsection}}
     {\newfloat{algorithm}{htbp}{loa}[subsection]}{}
  \ifthenelse{\equal{\ALG@within}{subsubsection}}
     {\newfloat{algorithm}{htbp}{loa}[subsubsection]}{}
  \ifthenelse{\equal{\ALG@within}{nothing}}
     {\newfloat{algorithm}{htbp}{loa}}{}
}{
  \newfloat{algorithm}{htbp}{loa}
}
\floatname{algorithm}{\ALG@name}

\newcommand{\listofalgorithms}{\listof{algorithm}{\listalgorithmname}}


\end{filecontents*}
\begin{filecontents*}[overwrite]{algorithmic.sty}
% ALGORITHMIC STYLE -- Released 8 APRIL 1996
%    for LaTeX version 2e
% Copyright -- 1994 Peter Williams
% E-mail PeterWilliams@dsto.defence.gov.au
%
% Modified by Alex Smola (08/2000)
% E-mail Alex.Smola@anu.edu.au
%
\NeedsTeXFormat{LaTeX2e}
\ProvidesPackage{algorithmic}
\typeout{Document Style `algorithmic' - environment}
%
\RequirePackage{ifthen}
\RequirePackage{calc}
\newboolean{ALC@noend}
\setboolean{ALC@noend}{false}
\newcounter{ALC@line}
\newcounter{ALC@rem}
\newlength{\ALC@tlm}
%
\DeclareOption{noend}{\setboolean{ALC@noend}{true}}
%
\ProcessOptions
%
% ALGORITHMIC
\newcommand{\algorithmicrequire}{\textbf{Require:}}
\newcommand{\algorithmicensure}{\textbf{Ensure:}}
\newcommand{\algorithmiccomment}[1]{\{#1\}}
\newcommand{\algorithmicend}{\textbf{end}}
\newcommand{\algorithmicif}{\textbf{if}}
\newcommand{\algorithmicthen}{\textbf{then}}
\newcommand{\algorithmicelse}{\textbf{else}}
\newcommand{\algorithmicelsif}{\algorithmicelse\ \algorithmicif}
\newcommand{\algorithmicendif}{\algorithmicend\ \algorithmicif}
\newcommand{\algorithmicfor}{\textbf{for}}
\newcommand{\algorithmicforall}{\textbf{for all}}
\newcommand{\algorithmicdo}{\textbf{do}}
\newcommand{\algorithmicendfor}{\algorithmicend\ \algorithmicfor}
\newcommand{\algorithmicwhile}{\textbf{while}}
\newcommand{\algorithmicendwhile}{\algorithmicend\ \algorithmicwhile}
\newcommand{\algorithmicloop}{\textbf{loop}}
\newcommand{\algorithmicendloop}{\algorithmicend\ \algorithmicloop}
\newcommand{\algorithmicrepeat}{\textbf{repeat}}
\newcommand{\algorithmicuntil}{\textbf{until}}

%changed by alex smola
\newcommand{\algorithmicinput}{\textbf{input}}
\newcommand{\algorithmicoutput}{\textbf{output}}
\newcommand{\algorithmicset}{\textbf{set}}
\newcommand{\algorithmictrue}{\textbf{true}}
\newcommand{\algorithmicfalse}{\textbf{false}}
\newcommand{\algorithmicand}{\textbf{and\ }}
\newcommand{\algorithmicor}{\textbf{or\ }}
\newcommand{\algorithmicfunction}{\textbf{function}}
\newcommand{\algorithmicendfunction}{\algorithmicend\ \algorithmicfunction}
\newcommand{\algorithmicmain}{\textbf{main}}
\newcommand{\algorithmicendmain}{\algorithmicend\ \algorithmicmain}
%end changed by alex smola

\def\ALC@item[#1]{%
\if@noparitem \@donoparitem
  \else \if@inlabel \indent \par \fi
         \ifhmode \unskip\unskip \par \fi
         \if@newlist \if@nobreak \@nbitem \else
                        \addpenalty\@beginparpenalty
                        \addvspace\@topsep \addvspace{-\parskip}\fi
           \else \addpenalty\@itempenalty \addvspace\itemsep
          \fi
    \global\@inlabeltrue
\fi
\everypar{\global\@minipagefalse\global\@newlistfalse
          \if@inlabel\global\@inlabelfalse \hskip -\parindent \box\@labels
             \penalty\z@ \fi
          \everypar{}}\global\@nobreakfalse
\if@noitemarg \@noitemargfalse \if@nmbrlist \refstepcounter{\@listctr}\fi \fi
\sbox\@tempboxa{\makelabel{#1}}%
\global\setbox\@labels
 \hbox{\unhbox\@labels \hskip \itemindent
       \hskip -\labelwidth \hskip -\ALC@tlm
       \ifdim \wd\@tempboxa >\labelwidth
                \box\@tempboxa
          \else \hbox to\labelwidth {\unhbox\@tempboxa}\fi
       \hskip \ALC@tlm}\ignorespaces}
%
\newenvironment{algorithmic}[1][0]{
\let\@item\ALC@item
  \newcommand{\ALC@lno}{%
\ifthenelse{\equal{\arabic{ALC@rem}}{0}}
{{\footnotesize \arabic{ALC@line}:}}{}%
}
\let\@listii\@listi
\let\@listiii\@listi
\let\@listiv\@listi
\let\@listv\@listi
\let\@listvi\@listi
\let\@listvii\@listi
  \newenvironment{ALC@g}{
    \begin{list}{\ALC@lno}{ \itemsep\z@ \itemindent\z@
    \listparindent\z@ \rightmargin\z@ 
    \topsep\z@ \partopsep\z@ \parskip\z@\parsep\z@
    \leftmargin 1em
    \addtolength{\ALC@tlm}{\leftmargin}
    }
  }
  {\end{list}}
  \newcommand{\ALC@it}{\addtocounter{ALC@line}{1}\addtocounter{ALC@rem}{1}\ifthenelse{\equal{\arabic{ALC@rem}}{#1}}{\setcounter{ALC@rem}{0}}{}\item}
  \newcommand{\ALC@com}[1]{\ifthenelse{\equal{##1}{default}}%
{}{\ \algorithmiccomment{##1}}}
  \newcommand{\REQUIRE}{\item[\algorithmicrequire]}
  \newcommand{\ENSURE}{\item[\algorithmicensure]}
  \newcommand{\STATE}{\ALC@it}
  \newcommand{\COMMENT}[1]{\algorithmiccomment{##1}}
%changes by alex smola
  \newcommand{\INPUT}{\item[\algorithmicinput]}
  \newcommand{\OUTPUT}{\item[\algorithmicoutput]}
  \newcommand{\SET}{\item[\algorithmicset]}
%  \newcommand{\TRUE}{\algorithmictrue}
%  \newcommand{\FALSE}{\algorithmicfalse}
  \newcommand{\AND}{\algorithmicand}
  \newcommand{\OR}{\algorithmicor}
  \newenvironment{ALC@func}{\begin{ALC@g}}{\end{ALC@g}}
  \newenvironment{ALC@main}{\begin{ALC@g}}{\end{ALC@g}}
%end changes by alex smola
  \newenvironment{ALC@if}{\begin{ALC@g}}{\end{ALC@g}}
  \newenvironment{ALC@for}{\begin{ALC@g}}{\end{ALC@g}}
  \newenvironment{ALC@whl}{\begin{ALC@g}}{\end{ALC@g}}
  \newenvironment{ALC@loop}{\begin{ALC@g}}{\end{ALC@g}}
  \newenvironment{ALC@rpt}{\begin{ALC@g}}{\end{ALC@g}}
  \renewcommand{\\}{\@centercr}
  \newcommand{\IF}[2][default]{\ALC@it\algorithmicif\ ##2\ \algorithmicthen%
\ALC@com{##1}\begin{ALC@if}}
  \newcommand{\SHORTIF}[2]{\ALC@it\algorithmicif\ ##1\
    \algorithmicthen\ {##2}}
  \newcommand{\ELSE}[1][default]{\end{ALC@if}\ALC@it\algorithmicelse%
\ALC@com{##1}\begin{ALC@if}}
  \newcommand{\ELSIF}[2][default]%
{\end{ALC@if}\ALC@it\algorithmicelsif\ ##2\ \algorithmicthen%
\ALC@com{##1}\begin{ALC@if}}
  \newcommand{\FOR}[2][default]{\ALC@it\algorithmicfor\ ##2\ \algorithmicdo%
\ALC@com{##1}\begin{ALC@for}}
  \newcommand{\FORALL}[2][default]{\ALC@it\algorithmicforall\ ##2\ %
\algorithmicdo%
\ALC@com{##1}\begin{ALC@for}}
  \newcommand{\SHORTFORALL}[2]{\ALC@it\algorithmicforall\ ##1\ %
    \algorithmicdo\ {##2}}
  \newcommand{\WHILE}[2][default]{\ALC@it\algorithmicwhile\ ##2\ %
\algorithmicdo%
\ALC@com{##1}\begin{ALC@whl}}
  \newcommand{\LOOP}[1][default]{\ALC@it\algorithmicloop%
\ALC@com{##1}\begin{ALC@loop}}
%changed by alex smola
  \newcommand{\FUNCTION}[2][default]{\ALC@it\algorithmicfunction\ ##2\ %
    \ALC@com{##1}\begin{ALC@func}}
  \newcommand{\MAIN}[2][default]{\ALC@it\algorithmicmain\ ##2\ %
    \ALC@com{##1}\begin{ALC@main}}
%end changed by alex smola
  \newcommand{\REPEAT}[1][default]{\ALC@it\algorithmicrepeat%
    \ALC@com{##1}\begin{ALC@rpt}}
    \newcommand{\UNTIL}[1]{\end{ALC@rpt}\ALC@it\algorithmicuntil\ ##1}
  \ifthenelse{\boolean{ALC@noend}}{
    \newcommand{\ENDIF}{\end{ALC@if}}
    \newcommand{\ENDFOR}{\end{ALC@for}}
    \newcommand{\ENDWHILE}{\end{ALC@whl}}
    \newcommand{\ENDLOOP}{\end{ALC@loop}}
    \newcommand{\ENDFUNCTION}{\end{ALC@func}}
    \newcommand{\ENDMAIN}{\end{ALC@main}}
  }{
    \newcommand{\ENDIF}{\end{ALC@if}\ALC@it\algorithmicendif}
    \newcommand{\ENDFOR}{\end{ALC@for}\ALC@it\algorithmicendfor}
    \newcommand{\ENDWHILE}{\end{ALC@whl}\ALC@it\algorithmicendwhile}
    \newcommand{\ENDLOOP}{\end{ALC@loop}\ALC@it\algorithmicendloop}
    \newcommand{\ENDFUNCTION}{\end{ALC@func}\ALC@it\algorithmicendfunction}
    \newcommand{\ENDMAIN}{\end{ALC@main}\ALC@it\algorithmicendmain}
  } 
  \renewcommand{\@toodeep}{}
  \begin{list}{\ALC@lno}{\setcounter{ALC@line}{0}\setcounter{ALC@rem}{0}%
      \itemsep\z@ \itemindent\z@ \listparindent\z@%
      \partopsep\z@ \parskip\z@ \parsep\z@%
      \labelsep 0.5em \topsep 0.2em%
      \ifthenelse{\equal{#1}{0}}
      {\labelwidth 0.5em }
      {\labelwidth  1.2em }
      \leftmargin\labelwidth \addtolength{\leftmargin}{\labelsep}
      \ALC@tlm\labelsep
      }
    }
  {\end{list}}















\end{filecontents*}
\begin{filecontents*}[overwrite]{fancyhdr.sty}
%%
%% This is file `fancyhdr.sty',
%% generated with the docstrip utility.
%%
%% The original source files were:
%%
%% fancyhdr.dtx  (with options: `fancyhdr')
%% 
%% This is a generated file.
%% 
%% This file may be distributed and/or modified under the conditions of
%% the LaTeX Project Public License, either version 1.3 of this license
%% or (at your option) any later version.  The latest version of this
%% license is in:
%% 
%%    http://www.latex-project.org/lppl.txt
%% 
%% and version 1.3 or later is part of all distributions of LaTeX version
%% 2005/12/01 or later.
%% 
%%%%%%%%%%%%%%%%%%%%%%%%%%%%%%%%%%%%%%%%%%%%%%%%%%%%%%%%%%%%%%%%%%%%%%%%
\NeedsTeXFormat{LaTeX2e}[2018-04-01]
\ProvidesPackage{fancyhdr}%
           [2025/02/07 v5.2
                  Extensive control of page headers and footers]%
% Copyright (C) 1994-2025 by Pieter van Oostrum <pieter@vanoostrum.org>
%%%%%%%%%%%%%%%%%%%%%%%%%%%%%%%%%%%%%%%%%%%%%%%%%%%%%%%%%%%%%%%%%%%%%%%%
\ifdefined\NewDocumentCommand\else\RequirePackage{xparse}\fi
\newif\iff@nch@check
\f@nch@checktrue
\DeclareOption{nocheck}{%
  \f@nch@checkfalse
}
\let\f@nch@gbl\relax
\newif\iff@nch@compatViii
\DeclareOption{compatV3}{%
  \PackageWarningNoLine{fancyhdr}{The `compatV3' option is deprecated.\MessageBreak
    It will disappear in one of the following releases.\MessageBreak
    Please change your document to work\MessageBreak
    without this option}
  \let\f@nch@gbl\global
  \f@nch@compatViiitrue
}
\newif\iff@nch@twoside
\f@nch@twosidefalse
\DeclareOption{twoside}{%
  \if@twoside\else\f@nch@twosidetrue\fi
}
\newcommand\f@nch@def[2]{%
  \def\temp@a{#2}\ifx\temp@a\@empty\f@nch@gbl\def#1{}%
                 \else\f@nch@gbl\def#1{#2\strut}\fi}
\DeclareOption{myheadings}{%
  \@ifundefined{chapter}{%
    \def\ps@myheadings{\ps@f@nch@fancyproto \let\@mkboth\@gobbletwo
      \fancyhf{}
      \fancyhead[LE,RO]{\thepage}%
      \fancyhead[RE]{\slshape\leftmark}%
      \fancyhead[LO]{\slshape\rightmark}%
      \let\sectionmark\@gobble
      \let\subsectionmark\@gobble
    }%
  }%
  {\def\ps@myheadings{\ps@f@nch@fancyproto \let\@mkboth\@gobbletwo
      \fancyhf{}
      \fancyhead[LE,RO]{\thepage}%
      \fancyhead[RE]{\slshape\leftmark}%
      \fancyhead[LO]{\slshape\rightmark}%
      \let\chaptermark\@gobble
      \let\sectionmark\@gobble
    }%
  }%
}
\DeclareOption{headings}{%
  \@ifundefined{chapter}{%
    \if@twoside
      \def\ps@headings{\ps@f@nch@fancyproto \def\@mkboth{\protect\markboth}
        \fancyhf{}
        \fancyhead[LE,RO]{\thepage}%
        \fancyhead[RE]{\slshape\leftmark}%
        \fancyhead[LO]{\slshape\rightmark}%
        \def\sectionmark##1{%
          \markboth{\MakeUppercase{%
            \ifnum \c@secnumdepth >\z@ \thesection\quad \fi##1}}{}}%
        \def\subsectionmark##1{%
          \markright{%
            \ifnum \c@secnumdepth >\@ne \thesubsection\quad \fi##1}}%
      }%
    \else
      \def\ps@headings{\ps@f@nch@fancyproto \def\@mkboth{\protect\markboth}
        \fancyhf{}
        \fancyhead[LE,RO]{\thepage}%
        \fancyhead[RE]{\slshape\leftmark}%
        \fancyhead[LO]{\slshape\rightmark}%
        \def\sectionmark##1{%
          \markright {\MakeUppercase{%
            \ifnum \c@secnumdepth >\z@ \thesection\quad \fi##1}}}%
        \let\subsectionmark\@gobble % Not needed but inserted for safety
      }%
    \fi
  }{\if@twoside
      \def\ps@headings{\ps@f@nch@fancyproto \def\@mkboth{\protect\markboth}
        \fancyhf{}
        \fancyhead[LE,RO]{\thepage}%
        \fancyhead[RE]{\slshape\leftmark}%
        \fancyhead[LO]{\slshape\rightmark}%
        \def\chaptermark##1{%
          \markboth{\MakeUppercase{%
            \ifnum \c@secnumdepth >\m@ne \if@mainmatter
              \@chapapp\ \thechapter. \ \fi\fi##1}}{}}%
        \def\sectionmark##1{%
          \markright {\MakeUppercase{%
            \ifnum \c@secnumdepth >\z@ \thesection. \ \fi##1}}}%
      }%
    \else
      \def\ps@headings{\ps@f@nch@fancyproto \def\@mkboth{\protect\markboth}
        \fancyhf{}
        \fancyhead[LE,RO]{\thepage}%
        \fancyhead[RE]{\slshape\leftmark}%
        \fancyhead[LO]{\slshape\rightmark}%
        \def\chaptermark##1{%
          \markright{\MakeUppercase{%
            \ifnum \c@secnumdepth >\m@ne \if@mainmatter
              \@chapapp\ \thechapter. \ \fi\fi##1}}}%
        \let\sectionmark\@gobble % Not needed but inserted for safety
      }%
    \fi
  }%
}
\ProcessOptions*
\newcommand{\f@nch@forc}[3]{\expandafter\f@nchf@rc\expandafter#1\expandafter{#2}{#3}}
\newcommand{\f@nchf@rc}[3]{\def\temp@ty{#2}\ifx\@empty\temp@ty\else
                                    \f@nch@rc#1#2\f@nch@rc{#3}\fi}
\long\def\f@nch@rc#1#2#3\f@nch@rc#4{\def#1{#2}#4\f@nchf@rc#1{#3}{#4}}
\newcommand{\f@nch@for}[3]{\edef\@fortmp{#2}%
  \expandafter\@forloop#2,\@nil,\@nil\@@#1{#3}}
\newcommand\f@nch@default[3]{%
  \edef\temp@a{\lowercase{\edef\noexpand\temp@a{#3}}}\temp@a \def#1{}%
  \f@nch@forc\tmpf@ra{#2}%
  {\expandafter\f@nch@ifin\tmpf@ra\temp@a{\edef#1{#1\tmpf@ra}}{}}%
  \ifx\@empty#1\def#1{#2}\fi}
\newcommand{\f@nch@ifin}[4]{%
  \edef\temp@a{#2}\def\temp@b##1#1##2\temp@b{\def\temp@b{##1}}%
  \expandafter\temp@b#2#1\temp@b\ifx\temp@a\temp@b #4\else #3\fi}
\newcommand{\fancyhead}[2][]{\f@nch@fancyhf\fancyhead h[#1]{#2}}%
\newcommand{\fancyfoot}[2][]{\f@nch@fancyhf\fancyfoot f[#1]{#2}}%
\newcommand{\fancyhf}[2][]{\f@nch@fancyhf\fancyhf {}[#1]{#2}}%
\newcommand{\fancyheadoffset}[2][]{\f@nch@fancyhfoffs\fancyheadoffset h[#1]{#2}}%
\newcommand{\fancyfootoffset}[2][]{\f@nch@fancyhfoffs\fancyfootoffset f[#1]{#2}}%
\newcommand{\fancyhfoffset}[2][]{\f@nch@fancyhfoffs\fancyhfoffset {}[#1]{#2}}%
\def\f@nch@fancyhf@Echeck#1{%
  \if@twoside\else
    \iff@nch@twoside\else
      \if\f@nch@@eo e%
        \PackageWarning{fancyhdr} {\string#1's `E' option without twoside option is useless.\MessageBreak
          Please consider using the `twoside' option}%
  \fi\fi\fi
}
\long\def\f@nch@fancyhf#1#2[#3]#4{%
  \def\temp@c{}%
  \f@nch@forc\tmpf@ra{#3}%
  {\expandafter\f@nch@ifin\tmpf@ra{eolcrhf,EOLCRHF}%
    {}{\edef\temp@c{\temp@c\tmpf@ra}}}%
  \ifx\@empty\temp@c\else \PackageError{fancyhdr}{Illegal char `\temp@c' in
    \string#1 argument: [#3]}{}%
  \fi \f@nch@for\temp@c{#3}%
  {\f@nch@default\f@nch@@eo{eo}\temp@c
    \f@nch@fancyhf@Echeck{#1}%
    \f@nch@default\f@nch@@lcr{lcr}\temp@c
    \f@nch@default\f@nch@@hf{hf}{#2\temp@c}%
    \f@nch@forc\f@nch@eo\f@nch@@eo
        {\f@nch@forc\f@nch@lcr\f@nch@@lcr
          {\f@nch@forc\f@nch@hf\f@nch@@hf
            {\expandafter\f@nch@def\csname
              f@nch@\f@nch@eo\f@nch@lcr\f@nch@hf\endcsname {#4}}}}}}
\def\f@nch@fancyhfoffs#1#2[#3]#4{%
  \def\temp@c{}%
  \f@nch@forc\tmpf@ra{#3}%
  {\expandafter\f@nch@ifin\tmpf@ra{eolrhf,EOLRHF}%
    {}{\edef\temp@c{\temp@c\tmpf@ra}}}%
  \ifx\@empty\temp@c\else \PackageError{fancyhdr}{Illegal char `\temp@c' in
    \string#1 argument: [#3]}{}%
  \fi \f@nch@for\temp@c{#3}%
  {\f@nch@default\f@nch@@eo{eo}\temp@c
    \f@nch@fancyhf@Echeck{#1}%
    \f@nch@default\f@nch@@lcr{lr}\temp@c
    \f@nch@default\f@nch@@hf{hf}{#2\temp@c}%
    \f@nch@forc\f@nch@eo\f@nch@@eo
        {\f@nch@forc\f@nch@lcr\f@nch@@lcr
          {\f@nch@forc\f@nch@hf\f@nch@@hf
            {\expandafter\setlength\csname
              f@nch@offset@\f@nch@eo\f@nch@lcr\f@nch@hf\endcsname {#4}}}}}%
  \f@nch@setoffs}
\NewDocumentCommand {\fancyheadwidth}{ s O{} O{} m }
                      {\f@nch@fancyhfwidth{#1}\fancyheadwidth h[#2][#3]{#4}}%
\NewDocumentCommand {\fancyfootwidth}{ s O{} O{} m }
                      {\f@nch@fancyhfwidth{#1}\fancyfootwidth f[#2][#3]{#4}}%
\NewDocumentCommand {\fancyhfwidth}  { s O{} O{} m }
                      {\f@nch@fancyhfwidth{#1}\fancyhfwidth  {}[#2][#3]{#4}}%
\def\f@nch@fancyhfwidth#1#2#3[#4][#5]#6{%
  \setlength\@tempdima{#6}%
  \def\temp@c{}%
  \f@nch@forc\tmpf@ra{#4}%
  {\expandafter\f@nch@ifin\tmpf@ra{eolcrhf,EOLCRHF}%
    {}{\edef\temp@c{\temp@c\tmpf@ra}}}%
  \ifx\@empty\temp@c\else \PackageError{fancyhdr}{Illegal char `\temp@c' in
    \string#2 argument: [#4]}{}%
  \fi
  \f@nch@for\temp@c{#4}%
  {\f@nch@default\f@nch@@eo{eo}\temp@c
    \f@nch@fancyhf@Echeck{#2}%
    \f@nch@default\f@nch@@lcr{lcr}\temp@c
    \f@nch@default\f@nch@@hf{hf}{#3\temp@c}%
    \f@nch@forc\f@nch@eo\f@nch@@eo
        {\f@nch@forc\f@nch@lcr\f@nch@@lcr
          {\f@nch@forc\f@nch@hf\f@nch@@hf
            {%
              \IfBooleanTF{#1}{%
                \expandafter\edef\csname
                  f@nch@width@\f@nch@eo\f@nch@lcr\f@nch@hf\endcsname{\the\@tempdima}%
              }%
              {%
                \expandafter\def\csname
                  f@nch@width@\f@nch@eo\f@nch@lcr\f@nch@hf\endcsname{#6}%
              }%
              \csname f@nchdrwdt@align@v@\f@nch@hf\endcsname
              \edef\f@nch@align@@h{\f@nch@lcr}%
              \def\temp@a{#5}%
              \ifx\temp@a\@empty \else \f@nchdrwdt@align#5\@nil{#2}\fi
              \expandafter\edef\csname
                f@nch@align@\f@nch@eo\f@nch@lcr\f@nch@hf\endcsname
                   {\f@nch@align@@v\f@nch@align@@h}}}}}}
\def\f@nch@width@elh{\headwidth}
\def\f@nch@width@ech{\headwidth}
\def\f@nch@width@erh{\headwidth}
\def\f@nch@width@olh{\headwidth}
\def\f@nch@width@och{\headwidth}
\def\f@nch@width@orh{\headwidth}
\def\f@nch@width@elf{\headwidth}
\def\f@nch@width@ecf{\headwidth}
\def\f@nch@width@erf{\headwidth}
\def\f@nch@width@olf{\headwidth}
\def\f@nch@width@ocf{\headwidth}
\def\f@nch@width@orf{\headwidth}
\def\f@nch@align@elh{bl}
\def\f@nch@align@ech{bc}
\def\f@nch@align@erh{br}
\def\f@nch@align@olh{bl}
\def\f@nch@align@och{bc}
\def\f@nch@align@orh{br}
\def\f@nch@align@elf{tl}
\def\f@nch@align@ecf{tc}
\def\f@nch@align@erf{tr}
\def\f@nch@align@olf{tl}
\def\f@nch@align@ocf{tc}
\def\f@nch@align@orf{tr}
\def\f@nchdrwdt@align@v@h{\def\f@nch@align@@v{b}}%
\def\f@nchdrwdt@align@v@f{\def\f@nch@align@@v{t}}%
\long\def\f@nchdrwdt@align#1#2\@nil#3{%
  \f@nch@ifin{#1}{TtcbB-}{%
    \f@nch@ifin{#1}{-}{}{\def\f@nch@align@@v{#1}}%
    \def\@tempa{#2}%
    \ifx\@tempa\@empty \else \def\f@nch@align@@h{#2}\fi
  }%
  {\def\f@nch@align@@h{#1}}%
  \expandafter\f@nch@ifin\expandafter{\f@nch@align@@h}{lcrj}{}%
    {\PackageError{fancyhdr}
                  {\string#3: Illegal char `\f@nch@align@@h'\MessageBreak
                              in alignment argument}{}}%
}
\newcommand{\lhead}[2][\f@nch@olh]%
                     {\f@nch@def\f@nch@olh{#2}\f@nch@def\f@nch@elh{#1}}
\newcommand{\chead}[2][\f@nch@och]%
                     {\f@nch@def\f@nch@och{#2}\f@nch@def\f@nch@ech{#1}}
\newcommand{\rhead}[2][\f@nch@orh]%
                     {\f@nch@def\f@nch@orh{#2}\f@nch@def\f@nch@erh{#1}}
\newcommand{\lfoot}[2][\f@nch@olf]%
                     {\f@nch@def\f@nch@olf{#2}\f@nch@def\f@nch@elf{#1}}
\newcommand{\cfoot}[2][\f@nch@ocf]%
                     {\f@nch@def\f@nch@ocf{#2}\f@nch@def\f@nch@ecf{#1}}
\newcommand{\rfoot}[2][\f@nch@orf]%
                     {\f@nch@def\f@nch@orf{#2}\f@nch@def\f@nch@erf{#1}}
\newlength{\f@nch@headwidth} \let\headwidth\f@nch@headwidth
\newlength{\f@nch@offset@elh}
\newlength{\f@nch@offset@erh}
\newlength{\f@nch@offset@olh}
\newlength{\f@nch@offset@orh}
\newlength{\f@nch@offset@elf}
\newlength{\f@nch@offset@erf}
\newlength{\f@nch@offset@olf}
\newlength{\f@nch@offset@orf}
\newcommand{\headrulewidth}{0.4pt}
\newcommand{\footrulewidth}{0pt}
\@ifundefined{headruleskip}%
      {\newcommand{\headruleskip}{0pt}}{}
\@ifundefined{footruleskip}%
      {\newcommand{\footruleskip}{.3\normalbaselineskip}}{}
\newcommand{\plainheadrulewidth}{0pt}
\newcommand{\plainfootrulewidth}{0pt}
\newif\if@fancyplain \@fancyplainfalse
\def\fancyplain#1#2{\if@fancyplain#1\else#2\fi}
\headwidth=-123456789sp
\let\f@nch@raggedleft\raggedleft
\let\f@nch@raggedright\raggedright
\let\f@nch@centering\centering
\let\f@nch@everypar\everypar
\ifdefined\ExplSyntaxOn
  \ExplSyntaxOn
  \providecommand\IfFormatAtLeastTF{\@ifl@t@r\fmtversion}
  \IfFormatAtLeastTF{2021-06-01}{
    \def\f@nch@saveclr@parhook #1{
      \expandafter\let\csname f@nch@__hook~#1\expandafter\endcsname
                      \csname __hook~#1\endcsname
      \expandafter\let\csname f@nch@__hook_toplevel~#1\expandafter\endcsname
                      \csname __hook_toplevel~#1\endcsname
      \expandafter\let\csname f@nch@__hook_next~#1\expandafter\endcsname
                      \csname __hook_next~#1\endcsname
      \expandafter\let\csname f@nch@g__hook_#1_code_prop\expandafter\endcsname
                      \csname g__hook_#1_code_prop\endcsname
      \RemoveFromHook{#1}[*]
      \ClearHookNext{#1}
    }
    \def\f@nch@restore@parhook #1{
      \global\expandafter\let\csname __hook~#1\expandafter\endcsname
                             \csname f@nch@__hook~#1\endcsname
      \global\expandafter\let\csname __hook_toplevel~#1\expandafter\endcsname
                             \csname f@nch@__hook_toplevel~#1\endcsname
      \global\expandafter\let\csname __hook_next~#1\expandafter\endcsname
                             \csname f@nch@__hook_next~#1\endcsname
      \global\expandafter\let\csname g__hook_#1_code_prop\expandafter\endcsname
                             \csname f@nch@g__hook_#1_code_prop\endcsname
    }
    \def\f@nch@resetpar{
      \f@nch@everypar{}
      \f@nch@saveclr@parhook{para/before}
      \f@nch@saveclr@parhook{para/begin}
      \f@nch@saveclr@parhook{para/end}
      \f@nch@saveclr@parhook{para/after}
    }
    \def\f@nch@restorepar{
      \f@nch@restore@parhook{para/before}
      \f@nch@restore@parhook{para/begin}
      \f@nch@restore@parhook{para/end}
      \f@nch@restore@parhook{para/after}
    }
  }{
    \def\f@nch@resetpar{
      \f@nch@everypar{}
    }
    \def\f@nch@restorepar{}
  }
  \ExplSyntaxOff
\else
  \def\f@nch@resetpar{%
    \f@nch@everypar{}%
  }
  \def\f@nch@restorepar{}
\fi
\newcommand\f@nch@noUppercase[2][]{#2}
\def\f@nch@reset{\f@nch@resetpar\restorecr\endlinechar=13
  \catcode`\\=0\catcode`\{=1\catcode`\}=2\catcode`\$=3\catcode`\&=4
  \catcode`\#=6\catcode`\^=7\catcode`\_=8\catcode`\ =10\catcode`\@=11
  \catcode`\:=11\catcode`\~=13\catcode`\%=14
  \catcode0=15 %NULL
  \catcode9=10 %TAB
  \let\\\@normalcr \let\raggedleft\f@nch@raggedleft
  \let\raggedright\f@nch@raggedright \let\centering\f@nch@centering
  \def\baselinestretch{1}%
  \hsize=\headwidth
  \def\nouppercase##1{{%
      \let\uppercase\relax\let\MakeUppercase\f@nch@noUppercase
      \expandafter\let\csname MakeUppercase \endcsname\relax
      \expandafter\def\csname MakeUppercase\space\space\space\endcsname
                                                   [####1]####2{####2}%
      ##1}}%
  \@ifundefined{@normalsize} {\normalsize} % for ucthesis.cls
   {\@normalsize}%
  }
\newcommand*{\fancycenter}[1][1em]{%
  \@ifnextchar[{\f@nch@center{#1}}{\f@nch@center{#1}[3]}%
}
\def\f@nch@center#1[#2]#3#4#5{%
  \def\@tempa{#4}\ifx\@tempa\@empty
    \hbox to\linewidth{\color@begingroup{#3}\hfil {#5}\color@endgroup}%
  \else
    \setlength\@tempdima{#1}%
    \setlength{\@tempdimb}{#2\@tempdima}%
    \@tempdimc \@tempdimb \advance\@tempdimc -\@tempdima
    \setlength\@tempskipa{\@tempdimb \@plus 1fil \@minus \@tempdimc}%
    \@tempskipb\@tempskipa
    \def\@tempa{#3}\ifx\@tempa\@empty
      \addtolength\@tempskipa{\z@ \@minus \@tempdima}%
    \fi
    \def\@tempa{#5}\ifx\@tempa\@empty % empty right
      \addtolength\@tempskipb{\z@ \@minus \@tempdima}%
    \fi
    \settowidth{\@tempdimb}{#3}%
    \settowidth{\@tempdimc}{#5}%
    \ifdim\@tempdimb>\@tempdimc
      \advance\@tempdimb -\@tempdimc
      \addtolength\@tempskipb{\@tempdimb \@minus \@tempdimb}%
    \else
      \advance\@tempdimc -\@tempdimb
      \addtolength\@tempskipa{\@tempdimc \@minus \@tempdimc}%
    \fi
    \hbox to\linewidth{\color@begingroup{#3}\hskip \@tempskipa
                      {#4}\hskip \@tempskipb {#5}\color@endgroup}%
  \fi
}
\newcommand{\f@nch@headinit}{}
\newcommand{\fancyheadinit}[1]{%
  \def\f@nch@headinit{#1}%
}
\newcommand{\f@nch@footinit}{}
\newcommand{\fancyfootinit}[1]{%
  \def\f@nch@footinit{#1}%
}
\newcommand{\fancyhfinit}[1]{%
  \def\f@nch@headinit{#1}%
  \def\f@nch@footinit{#1}%
}
\ifdefined\NewMirroredHookPair
  \NewMirroredHookPair{fancyhdr/before}{fancyhdr/after}
  \NewMirroredHookPair{fancyhdr/head/begin}{fancyhdr/head/end}
  \NewMirroredHookPair{fancyhdr/foot/begin}{fancyhdr/foot/end}
\fi
\newlength\f@nch@height
\newlength\f@nch@footalignment
\newif\iff@nch@footalign\f@nch@footalignfalse
\newcommand{\fancyfootalign}[1]{%
  \def\temp@a{#1}%
  \ifx\temp@a\@empty
    \f@nch@footalignfalse
  \else
    \f@nch@footaligntrue
    \setlength\f@nch@footalignment{#1}%
  \fi
}
\newcommand\fancyhdrsettoheight[2]{%
  \expandafter\ifx\csname f@nch@#2\endcsname\fancyhdrsettoheight
    \else\PackageError{fancyhdr}{Unknown parameter #2 in \string\fancyhdrsettoheight}{}\fi
  \setbox\@tempboxa\hbox{{\f@nch@checkfalse\csname @#2\endcsname}}%
  \setlength{#1}\f@nch@height
  \setbox\@tempboxa\box\voidb@x
}
\let\f@nch@oddhead\fancyhdrsettoheight
\let\f@nch@evenhead\fancyhdrsettoheight
\let\f@nch@oddfoot\fancyhdrsettoheight
\let\f@nch@evenfoot\fancyhdrsettoheight
\newcommand\f@nch@vbox[2]{%
  \setbox0\vbox{#2}%
  \global\f@nch@height=\ht0
  \ifdim\ht0>#1\relax
    \iff@nch@check
      \dimen0=#1\advance\dimen0-\ht0
      \PackageWarning{fancyhdr}{%
        \string#1 is too small (\the#1): \MessageBreak
        Make it at least \the\ht0, for example:\MessageBreak
        \string\setlength{\string#1}{\the\ht0}%
        \iff@nch@compatViii .\MessageBreak
        We now make it that large for the rest of the document.\MessageBreak
        This may cause the page layout to be inconsistent, however
        \fi
        \ifx#1\headheight .\MessageBreak
          You might also make \topmargin smaller:\MessageBreak
          \string\addtolength{\string\topmargin}{\the\dimen0}%
        \fi
        \@gobble
      }%
      \iff@nch@compatViii
        \dimen0=#1\relax
        \global#1=\ht0\relax
        \ht0=\dimen0 %
      \else
        \ht0=#1\relax
      \fi
    \else
      \ht0=#1\relax
    \fi
  \fi
  \box0}
\newcommand\f@nch@head[6]{%
  \f@nch@reset
  \ifdefined\UseHook\UseHook{fancyhdr/before}\UseHook{fancyhdr/head/begin}\fi
  \f@nch@headinit\relax
  #1%
  \hbox to\headwidth{%
    \f@nch@vbox\headheight{%
      \f@nch@hfbox{#2}{#3}{#4}{#6}{h}%
      \vskip\headruleskip\relax
      \headrule
    }%
  }%
  #5%
  \ifdefined\UseHook\UseHook{fancyhdr/head/end}\UseHook{fancyhdr/after}\fi
  \f@nch@restorepar
}
\newcommand\f@nch@foot[6]{%
  \f@nch@reset
  \ifdefined\UseHook\UseHook{fancyhdr/before}\UseHook{fancyhdr/foot/begin}\fi
  \f@nch@footinit\relax
  #1%
  \hbox to\headwidth{%
    \f@nch@vbox\footskip{%
      \setbox0=\vbox{\footrule}\unvbox0
      \vskip\footruleskip
      \f@nch@hfbox{#2}{#3}{#4}{#6}{f}%
    \iff@nch@footalign \vskip\f@nch@footalignment \fi
    }%
  }%
  #5%
  \ifdefined\UseHook\UseHook{fancyhdr/foot/end}\UseHook{fancyhdr/after}\fi
  \f@nch@restorepar
}
\newlength\f@nch@widthL
\newlength\f@nch@widthC
\newlength\f@nch@widthR
\newcommand\f@nch@hfbox[5]{%
  \setlength\f@nch@widthL{\csname f@nch@width@#4l#5\endcsname}%
  \setlength\f@nch@widthC{\csname f@nch@width@#4c#5\endcsname}%
  \setlength\f@nch@widthR{\csname f@nch@width@#4r#5\endcsname}%
  \let\@tempa\f@nch@hfbox@center
  \ifdim \dimexpr \f@nch@widthL+\f@nch@widthC+\f@nch@widthR>\headwidth
  \else
    \ifdim \dimexpr \f@nch@widthL+0.5\f@nch@widthC>0.5\headwidth
      \let \@tempa\f@nch@hfbox@fit
    \fi
    \ifdim \dimexpr \f@nch@widthR+0.5\f@nch@widthC>0.5\headwidth
      \let \@tempa\f@nch@hfbox@fit
    \fi
  \fi
  \@tempa{#1}{#2}{#3}#4#5%
}
\newcommand\f@nch@hfbox@center[5]{%
  \hbox to \headwidth{%
    \rlap{\f@nch@parbox{#1}\f@nch@widthL{#4}l{#5}}%
    \hfill
    \f@nch@parbox{#2}\f@nch@widthC{#4}c{#5}%
    \hfill
    \llap{\f@nch@parbox{#3}\f@nch@widthR{#4}r{#5}}%
  }%
}
\newcommand\f@nch@hfbox@fit[5]{%
  \hbox to \headwidth{%
    \f@nch@parbox{#1}\f@nch@widthL{#4}l{#5}%
    \hfill
    \f@nch@parbox{#2}\f@nch@widthC{#4}c{#5}%
    \hfill
    \f@nch@parbox{#3}\f@nch@widthR{#4}r{#5}%
  }%
}%
\newcommand\f@nch@parbox[5]{%
  \expandafter\expandafter\expandafter\f@nch@parbox@align
                     \csname f@nch@align@#3#4#5\endcsname
  \parbox[\f@nch@align@@v]{#2}%
    {%
      \f@nch@align@@pre
      \f@nch@align@@h\leavevmode\ignorespaces#1%
      \f@nch@align@@post
    }%
}
\newcommand\f@nch@parbox@align[2]{%
  \def\f@nch@align@@pre{}%
  \def\f@nch@align@@post{}%
  \csname f@nch@parbox@align@v#1\endcsname
  \csname f@nch@parbox@align@h#2\endcsname
}
\def\f@nch@parbox@align@vT{\def\f@nch@align@@v{t}\def\f@nch@align@@pre{\vspace{0pt}}}
\def\f@nch@parbox@align@vt{\def\f@nch@align@@v{t}}
\def\f@nch@parbox@align@vc{\def\f@nch@align@@v{c}}
\def\f@nch@parbox@align@vb{\def\f@nch@align@@v{b}}
\def\f@nch@parbox@align@vB{\def\f@nch@align@@v{b}\def\f@nch@align@@post{\vspace{0pt}}}
\def\f@nch@parbox@align@hl{\def\f@nch@align@@h{\raggedright}}
\def\f@nch@parbox@align@hc{\def\f@nch@align@@h{\centering}}
\def\f@nch@parbox@align@hr{\def\f@nch@align@@h{\raggedleft}}
\def\f@nch@parbox@align@hj{\def\f@nch@align@@h{}}
\@ifundefined{@chapapp}{\let\@chapapp\chaptername}{}%
\def\f@nch@initialise{%
  \@ifundefined{chapter}%
   {\def\sectionmark##1{\markboth{\MakeUppercase{\ifnum \c@secnumdepth>\z@
          \thesection\hskip 1em\relax
        \fi ##1}}{}}%
    \def\subsectionmark##1{\markright {\ifnum \c@secnumdepth >\@ne
      \thesubsection\hskip 1em\relax \fi ##1}}}%
   {\def\chaptermark##1{\markboth {\MakeUppercase{\ifnum
        \c@secnumdepth>\m@ne \@chapapp\ \thechapter. \ \fi ##1}}{}}%
    \def\sectionmark##1{\markright{\MakeUppercase{\ifnum \c@secnumdepth >\z@
        \thesection. \ \fi ##1}}}%
   }%
  \def\headrule{{\if@fancyplain\let\headrulewidth\plainheadrulewidth\fi
      \hrule\@height\headrulewidth\@width\headwidth
      \vskip-\headrulewidth}}%
  \def\footrule{{\if@fancyplain\let\footrulewidth\plainfootrulewidth\fi
      \hrule\@width\headwidth\@height\footrulewidth}}%
  \def\headrulewidth{0.4pt}%
  \def\footrulewidth{0pt}%
  \def\headruleskip{0pt}%
  \def\footruleskip{0.3\normalbaselineskip}%
  \fancyhf{}%
  \if@twoside
    \fancyhead[el,or]{\fancyplain{}{\slshape\rightmark}}%
    \fancyhead[er,ol]{\fancyplain{}{\slshape\leftmark}}%
  \else
    \fancyhead[l]{\fancyplain{}{\slshape\rightmark}}%
    \fancyhead[r]{\fancyplain{}{\slshape\leftmark}}%
  \fi
  \fancyfoot[c]{\rmfamily\thepage}% page number
}
\f@nch@initialise
\def\ps@f@nch@fancyproto{%
  \ifdim\headwidth<0sp
    \global\advance\headwidth123456789sp\global\advance\headwidth\textwidth
  \fi
  \gdef\ps@f@nch@fancyproto{\@fancyplainfalse\ps@f@nch@fancycore}%
  \@fancyplainfalse\ps@f@nch@fancycore
}%
\@namedef{f@nch@ps@f@nch@fancyproto-is-fancyhdr}{}
\def\ps@fancy{\ps@f@nch@fancyproto}
\@namedef{f@nch@ps@fancy-is-fancyhdr}{}
\def\ps@fancyplain{\ps@f@nch@fancyproto \let\ps@plain\ps@plain@fancy}
\def\ps@plain@fancy{\@fancyplaintrue\ps@f@nch@fancycore}
\let\f@nch@ps@empty\ps@empty
\def\ps@f@nch@fancycore{%
  \f@nch@ps@empty
  \def\@mkboth{\protect\markboth}%
  \def\f@nch@oddhead{\f@nch@head\f@nch@Oolh\f@nch@olh\f@nch@och\f@nch@orh\f@nch@Oorh{o}}%
  \def\@oddhead{%
    \iff@nch@twoside
      \ifodd\c@page
        \f@nch@oddhead
      \else
        \@evenhead
      \fi
    \else
      \f@nch@oddhead
    \fi
  }
  \def\f@nch@oddfoot{\f@nch@foot\f@nch@Oolf\f@nch@olf\f@nch@ocf\f@nch@orf\f@nch@Oorf{o}}%
  \def\@oddfoot{%
    \iff@nch@twoside
      \ifodd\c@page
        \f@nch@oddfoot
      \else
        \@evenfoot
      \fi
    \else
      \f@nch@oddfoot
    \fi
  }
  \def\@evenhead{\f@nch@head\f@nch@Oelh\f@nch@elh\f@nch@ech\f@nch@erh\f@nch@Oerh{e}}%
  \def\@evenfoot{\f@nch@foot\f@nch@Oelf\f@nch@elf\f@nch@ecf\f@nch@erf\f@nch@Oerf{e}}%
}
\def\f@nch@Oolh{\if@reversemargin\hss\else\relax\fi}
\def\f@nch@Oorh{\if@reversemargin\relax\else\hss\fi}
\let\f@nch@Oelh\f@nch@Oorh
\let\f@nch@Oerh\f@nch@Oolh
\let\f@nch@Oolf\f@nch@Oolh
\let\f@nch@Oorf\f@nch@Oorh
\let\f@nch@Oelf\f@nch@Oelh
\let\f@nch@Oerf\f@nch@Oerh
\def\f@nch@offsolh{\headwidth=\textwidth\advance\headwidth\f@nch@offset@olh
                   \advance\headwidth\f@nch@offset@orh\hskip-\f@nch@offset@olh}
\def\f@nch@offselh{\headwidth=\textwidth\advance\headwidth\f@nch@offset@elh
                   \advance\headwidth\f@nch@offset@erh\hskip-\f@nch@offset@elh}
\def\f@nch@offsolf{\headwidth=\textwidth\advance\headwidth\f@nch@offset@olf
                   \advance\headwidth\f@nch@offset@orf\hskip-\f@nch@offset@olf}
\def\f@nch@offself{\headwidth=\textwidth\advance\headwidth\f@nch@offset@elf
                   \advance\headwidth\f@nch@offset@erf\hskip-\f@nch@offset@elf}
\def\f@nch@setoffs{%
  \f@nch@gbl\let\headwidth\f@nch@headwidth
  \f@nch@gbl\def\f@nch@Oolh{\f@nch@offsolh}%
  \f@nch@gbl\def\f@nch@Oelh{\f@nch@offselh}%
  \f@nch@gbl\def\f@nch@Oorh{\hss}%
  \f@nch@gbl\def\f@nch@Oerh{\hss}%
  \f@nch@gbl\def\f@nch@Oolf{\f@nch@offsolf}%
  \f@nch@gbl\def\f@nch@Oelf{\f@nch@offself}%
  \f@nch@gbl\def\f@nch@Oorf{\hss}%
  \f@nch@gbl\def\f@nch@Oerf{\hss}%
}
\newif\iff@nch@footnote
\AtBeginDocument{%
  \let\latex@makecol\@makecol
  \def\@makecol{\ifvoid\footins\f@nch@footnotefalse\else\f@nch@footnotetrue\fi
    \let\f@nch@topfloat\@toplist\let\f@nch@botfloat\@botlist\latex@makecol}%
}
\newcommand\iftopfloat[2]{\ifx\f@nch@topfloat\@empty #2\else #1\fi}%
\newcommand\ifbotfloat[2]{\ifx\f@nch@botfloat\@empty #2\else #1\fi}%
\newcommand\iffloatpage[2]{\if@fcolmade #1\else #2\fi}%
\newcommand\iffootnote[2]{\iff@nch@footnote #1\else #2\fi}%
\ifx\@temptokenb\undefined \csname newtoks\endcsname\@temptokenb\fi
\newif\iff@nch@pagestyle@star
\newcommand\fancypagestyle{%
  \@ifstar{\f@nch@pagestyle@startrue\f@nch@pagestyle}%
          {\f@nch@pagestyle@starfalse\f@nch@pagestyle}%
}
\newcommand\f@nch@pagestyle[1]{%
  \@ifnextchar[{\f@nch@@pagestyle{#1}}{\f@nch@@pagestyle{#1}[f@nch@fancyproto]}%
}
\long\def\f@nch@@pagestyle#1[#2]#3{%
  \@ifundefined{ps@#2}{%
    \PackageError{fancyhdr}{\string\fancypagestyle: Unknown base page style `#2'}{}%
  }{%
    \@ifundefined{f@nch@ps@#2-is-fancyhdr}{%
      \PackageError{fancyhdr}{\string\fancypagestyle: Base page style `#2' is not fancyhdr-based}{}%
    }%
    {%
      \f@nch@pagestyle@setup
      \def\temp@b{\@namedef{ps@#1}}%
      \expandafter\temp@b\expandafter{\the\@temptokenb
          \let\f@nch@gbl\relax\@nameuse{ps@#2}#3\relax}%
      \@namedef{f@nch@ps@#1-is-fancyhdr}{}%
    }%
  }%
}
\newcommand\f@nch@pagestyle@setup{%
  \iff@nch@pagestyle@star
    \iff@nch@check\@temptokenb={\f@nch@checktrue}\else\@temptokenb={\f@nch@checkfalse}\fi
    \@tfor\temp@a:=
      \f@nch@olh\f@nch@och\f@nch@orh\f@nch@elh\f@nch@ech\f@nch@erh
      \f@nch@olf\f@nch@ocf\f@nch@orf\f@nch@elf\f@nch@ecf\f@nch@erf
      \f@nch@width@elh\f@nch@width@ech\f@nch@width@erh\f@nch@width@olh
      \f@nch@width@och\f@nch@width@orh\f@nch@width@elf\f@nch@width@ecf
      \f@nch@width@erf\f@nch@width@olf\f@nch@width@ocf\f@nch@width@orf
      \f@nch@align@elh\f@nch@align@ech\f@nch@align@erh\f@nch@align@olh
      \f@nch@align@och\f@nch@align@orh\f@nch@align@elf\f@nch@align@ecf
      \f@nch@align@erf\f@nch@align@olf\f@nch@align@ocf\f@nch@align@orf
      \f@nch@Oolh\f@nch@Oorh\f@nch@Oelh\f@nch@Oerh
      \f@nch@Oolf\f@nch@Oorf\f@nch@Oelf\f@nch@Oerf
      \f@nch@headinit\f@nch@footinit
      \headrule\headrulewidth\footrule\footrulewidth
    \do {%
      \toks@=\expandafter\expandafter\expandafter{\temp@a}%
      \toks@=\expandafter\expandafter\expandafter{%
        \expandafter\expandafter\expandafter\def
        \expandafter\expandafter\temp@a\expandafter{\the\toks@}}%
      \edef\temp@b{\@temptokenb={\the\@temptokenb\the\toks@}}%
      \temp@b
    }%
    \@tfor\temp@a:=
      \f@nch@offset@olh\f@nch@offset@orh\f@nch@offset@elh\f@nch@offset@erh
      \f@nch@offset@olf\f@nch@offset@orf\f@nch@offset@elf\f@nch@offset@erf
    \do {%
      \toks@=\expandafter\expandafter\expandafter{\expandafter\the\temp@a}%
      \toks@=\expandafter\expandafter\expandafter{%
        \expandafter\expandafter\expandafter\setlength
        \expandafter\expandafter\temp@a\expandafter{\the\toks@}}%
      \edef\temp@b{\@temptokenb={\the\@temptokenb\the\toks@}}%
      \temp@b
    }%
  \else
    \@temptokenb={}%
  \fi
}
\newcommand\fancypagestyleassign[2]{%
  \@ifundefined{ps@#2}{%
    \PackageError{fancyhdr}{\string\fancypagestyleassign: Unknown page style `#2'}{}%
    }{%
     \expandafter\let
       \csname ps@#1\expandafter\endcsname
       \csname ps@#2\endcsname
     \@ifundefined{f@nch@ps@#2-is-fancyhdr}{%
       \expandafter\let\csname f@nch@ps@#1-is-fancyhdr\endcsname\@undefined
     }{%
       \@namedef{f@nch@ps@#1-is-fancyhdr}{}%
     }%
   }%
}
\fancypagestyle*{fancydefault}{\f@nch@initialise}
\def\f@nchdrbox@topstrut{\vrule height\ht\strutbox width\z@}
\def\f@nchdrbox@botstrut{\vrule depth\dp\strutbox width\z@}
\def\f@nchdrbox@nostrut{\noalign{\vspace{0pt}}\let\f@nchdrbox@@crstrut\f@nchdrbox@botstrut}
\NewDocumentCommand{\fancyhdrbox}{ O{cl} o m }{%
\begingroup
  \let\f@nchdrbox@@pre\f@nchdrbox@topstrut
  \let\f@nchdrbox@@postx\f@nchdrbox@botstrut
  \let\f@nchdrbox@@posty\relax
  \let\f@nchdrbox@@crstrut\strut
  \IfNoValueTF{#2}%
    {\let\f@nchdrbox@@halignto\@empty}%
    {\setlength\@tempdima{#2}%
      \def\f@nchdrbox@@halignto{to\@tempdima}}%
  \def\@tempa{#1}%
  \ifx\@tempa\@empty
    \f@nchdrbox@align cl\@nil{#3}%
  \else
    \f@nchdrbox@align #1\@nil{#3}%
  \fi
\endgroup
}
\protected\def\f@nchdrbox@cr{%
  {\ifnum0=`}\fi\@ifstar\@f@nchdrbox@xcr\@f@nchdrbox@xcr}

\def\@f@nchdrbox@xcr{%
  \unskip\f@nchdrbox@@crstrut
  \@ifnextchar[\@f@nchdrbox@argc{\ifnum0=`{\fi}\cr}%
}

\def\@f@nchdrbox@argc[#1]{%
  \ifnum0=`{\fi}%
    \ifdim #1>\z@
      \unskip\@f@nchdrbox@xargc{#1}%
    \else
      \@f@nchdrbox@yargc{#1}%
    \fi}

\def\@f@nchdrbox@xargc#1{\@tempdima #1\advance\@tempdima \dp \strutbox
   \vrule \@height\z@ \@depth\@tempdima \@width\z@ \cr}

\def\@f@nchdrbox@yargc#1{\cr\noalign{\setlength\@tempdima{#1}\vskip\@tempdima}}
\def\f@nchdrbox@T{\let\f@nchdrbox@@pre\f@nchdrbox@nostrut
                  \f@nchdrbox@t}
\def\f@nchdrbox@t{\def\f@nchdrbox@@v{t}\def\f@nchdrbox@@h{l}}
\def\f@nchdrbox@c{\def\f@nchdrbox@@v{c}\def\f@nchdrbox@@h{c}}
\def\f@nchdrbox@b{\def\f@nchdrbox@@v{b}\def\f@nchdrbox@@h{l}}
\def\f@nchdrbox@B{\let\f@nchdrbox@@postx\relax
                  \def\f@nchdrbox@@posty{\vspace{0pt}}%
                  \f@nchdrbox@b}
\long\def\f@nchdrbox@align#1#2\@nil#3{%
  \f@nch@ifin{#1}{TtcbB}{%
    \@nameuse{f@nchdrbox@#1}%
    \def\@tempa{#2}%
    \ifx\@tempa\@empty\else \def\f@nchdrbox@@h{#2}\fi
  }%
  {\def\f@nchdrbox@@v{c}\def\f@nchdrbox@@h{#1}}%
  \expandafter\f@nch@ifin\expandafter{\f@nchdrbox@@h}{lcr}{}%
  {\PackageError{fancyhdr}{\string\fancyhdrbox: Illegal char `\f@nchdrbox@@h'\MessageBreak
                            in alignment argument}{}}%
  \let\\\f@nchdrbox@cr
  \setbox0=\if \f@nchdrbox@@v t\vtop
  \else \vbox
  \fi
  {%
     \ialign \f@nchdrbox@@halignto
     \bgroup \relax
     {\if \f@nchdrbox@@h l\hskip 1sp\else \hfil \fi
       \ignorespaces ##\unskip
       \if\f@nchdrbox@@h r\else \hfil \fi
     }%
     \tabskip\z@skip \cr
     \f@nchdrbox@@pre
     #3\unskip \f@nchdrbox@@postx
     \crcr
     \egroup
     \f@nchdrbox@@posty
  }%
  \if\f@nchdrbox@@v c\@tempdima=\ht0\advance\@tempdima\dp0%
    \ht0=0.5\@tempdima\dp0=0.5\@tempdima\fi
  \leavevmode \box0
}
\@ifclassloaded{newlfm}
{
  \let\ps@@empty\f@nch@ps@empty
  \AtBeginDocument{%
    \renewcommand{\@zfancyhead}[5]{\relax\hbox to\headwidth{\f@nch@reset
      \@zfancyvbox\headheight{\hbox
        {\rlap{\parbox[b]{\headwidth}{\raggedright\f@nch@olh}}\hfill
          \parbox[b]{\headwidth}{\centering\f@nch@olh}\hfill
          \llap{\parbox[b]{\headwidth}{\raggedleft\f@nch@orh}}}%
        \zheadrule}}\relax}%
  }
}
{}
\endinput
%%
%% End of file `fancyhdr.sty'.

\end{filecontents*}
\begin{filecontents*}[overwrite]{icml2026.sty}
% File: icml2026.sty (LaTeX style file for ICML-2026, version of 2025-10-29)

% This file contains the LaTeX formatting parameters for a two-column
% conference proceedings that is 8.5 inches wide by 11 inches high.
%
% Modified by Hanze Dong, Alberto Bietti, and Felix Berkenkamp, 2025
% - Revert to times for better compatibility
% - Updated years, volume, location
% - Added preprint version
% - Based on the suggestion from Johan Larsson:
%   1. Added an end-of-document safety check to ensure the affiliations or notice footnote is printed:
%      (1) Introduces a flag \newif\ificml@noticeprinted and sets it false by default.
%      (2) At end of document, emits a package warning if \printAffiliationsAndNotice{...} was never called.
%   2. \printAffiliationsAndNotice now sets the flag when called: Begins with \global\icml@noticeprintedtrue.
% - Migrated to more recent version of fancyhdr for running title in header
%
% Modified by Johan Larsson, 2025
% - Use newtx instead of times, aligning serif, sans-serif, typerwriter,
%   and math fonts.
% - Use caption package to setup captions instead of manually defining themanually defining them.
% - Formatted icml2026.sty and example_paper.tex
% - Use title case for section title to 2.9
% - Replace subfigure package with subcaption in example, since it is
%   designed to work together with the caption package (which is now required).
% - Remove unused label in example
%
% Modified by Tegan Maharaj and Felix Berkenkamp 2025: changed years, volume, location
%
% Modified by Jonathan Scarlett 2024: changed years, volume, location
%
% Modified by Sivan Sabato 2023: changed years and volume number.
% Modified by Jonathan Scarlett 2023: added page numbers to every page
%
% Modified by Csaba Szepesvari 2022: changed years, PMLR ref. Turned off checking marginparwidth
%     as marginparwidth only controls the space available for margin notes and margin notes
%     will NEVER be used anyways in submitted versions, so there is no reason one should
%     check whether marginparwidth has been tampered with.
%     Also removed pdfview=FitH from hypersetup as it did not do its job; the default choice is a bit better
%     but of course the double-column format is not supported by this hyperlink preview functionality
%     in a completely satisfactory fashion.
% Modified by Gang Niu 2022: Changed color to xcolor
%
% Modified by Iain Murray 2018: changed years, location. Remove affiliation notes when anonymous.
%     Move times dependency from .tex to .sty so fewer people delete it.
%
% Modified by Daniel Roy 2017: changed byline to use footnotes for affiliations, and removed emails
%
% Modified by Percy Liang 12/2/2013: changed the year, location from the previous template for ICML 2014

% Modified by Fei Sha 9/2/2013: changed the year, location form the previous template for ICML 2013
%
% Modified by Fei Sha 4/24/2013: (1) remove the extra whitespace after the
%     first author's email address (in %the camera-ready version) (2) change the
%     Proceeding ... of ICML 2010 to 2014 so PDF's metadata will show up %
%     correctly
%
% Modified by Sanjoy Dasgupta, 2013: changed years, location
%
% Modified by Francesco Figari, 2012: changed years, location
%
% Modified by Christoph Sawade and Tobias Scheffer, 2011: added line
% numbers, changed years
%
% Modified by Hal Daume III, 2010: changed years, added hyperlinks
%
% Modified by Kiri Wagstaff, 2009: changed years
%
% Modified by Sam Roweis, 2008: changed years
%
% Modified by Ricardo Silva, 2007: update of the ifpdf verification
%
% Modified by Prasad Tadepalli and Andrew Moore, merely changing years.
%
% Modified by Kristian Kersting, 2005, based on Jennifer Dy's 2004 version
% - running title. If the original title is to long or is breaking a line,
%   use \icmltitlerunning{...} in the preamble to supply a shorter form.
%   Added fancyhdr package to get a running head.
% - Updated to store the page size because pdflatex does compile the
%   page size into the pdf.
%
% Hacked by Terran Lane, 2003:
% - Updated to use LaTeX2e style file conventions (ProvidesPackage,
%   etc.)
% - Added an ``appearing in'' block at the base of the first column
%   (thus keeping the ``appearing in'' note out of the bottom margin
%   where the printer should strip in the page numbers).
% - Added a package option [accepted] that selects between the ``Under
%   review'' notice (default, when no option is specified) and the
%   ``Appearing in'' notice (for use when the paper has been accepted
%   and will appear).
%
%   Originally created as:  ml2k.sty (LaTeX style file for ICML-2000)
%   by P. Langley (12/23/99)

%%%%%%%%%%%%%%%%%%%%
%% This version of the style file supports both a ``review'' version
%% and a ``final/accepted'' version.  The difference is only in the
%% text that appears in the note at the bottom of the first column of
%% the first page.  The default behavior is to print a note to the
%% effect that the paper is under review and don't distribute it.  The
%% final/accepted version prints an ``Appearing in'' note.  To get the
%% latter behavior, in the calling file change the ``usepackage'' line
%% from:
%% \usepackage{icml2025}
%% to
%% \usepackage[accepted]{icml2025}
%%%%%%%%%%%%%%%%%%%%

\NeedsTeXFormat{LaTeX2e}
\ProvidesPackage{icml2026}[2025/10/29 v2.0 ICML Conference Style File]

% Before 2018, \usepackage{times} was in the example TeX, but inevitably
% not everybody did it.
% \RequirePackage[amsthm]{newtx}
% 2025.11.6 revert to times for better compatibility
\RequirePackage{times}

% Use fancyhdr package
\RequirePackage{fancyhdr}
\RequirePackage{xcolor} % changed from color to xcolor (2021/11/24)
\RequirePackage{algorithm}
\RequirePackage{algorithmic}
\RequirePackage{natbib}
\RequirePackage{eso-pic} % used by \AddToShipoutPicture
\RequirePackage{forloop}
\RequirePackage{url}
\RequirePackage{caption}

%%%%%%%% Options
\DeclareOption{accepted}{%
  \renewcommand{\Notice@String}{\ICML@appearing}
  \gdef\isaccepted{1}
}

% === Preprint option ===
\DeclareOption{preprint}{%%
  \renewcommand{\Notice@String}{\ICML@preprint}%%
  \gdef\ispreprint{1}%%
}

% Distinct preprint footer text
\newcommand{\ICML@preprint}{%
  \textit{Preprint. \today.}%
}

\DeclareOption{nohyperref}{%
  \gdef\nohyperref{1}
}

% Helper flag: show real authors for accepted or preprint
\newif\ificmlshowauthors
\icmlshowauthorsfalse

%%%%%%%%%%%%%%%%%%%%
% This string is printed at the bottom of the page for the
% final/accepted version of the ``appearing in'' note.  Modify it to
% change that text.
%%%%%%%%%%%%%%%%%%%%
\newcommand{\ICML@appearing}{\textit{Proceedings of the
$\mathit{43}^{rd}$ International Conference on Machine Learning},
Seoul, South Korea. PMLR 306, 2026.
Copyright 2026 by the author(s).}

%%%%%%%%%%%%%%%%%%%%
% This string is printed at the bottom of the page for the draft/under
% review version of the ``appearing in'' note.  Modify it to change
% that text.
%%%%%%%%%%%%%%%%%%%%
\newcommand{\Notice@String}{Preliminary work.  Under review by the
International Conference on Machine Learning (ICML)\@.  Do not distribute.}

% Cause the declared options to actually be parsed and activated
\ProcessOptions\relax

% After options are processed, decide if authors should be visible
\ifdefined\isaccepted \icmlshowauthorstrue \fi
\ifdefined\ispreprint \icmlshowauthorstrue \fi

\ifdefined\isaccepted\else\ifdefined\ispreprint\else\ifdefined\hypersetup
  \hypersetup{pdfauthor={Anonymous Authors}}
\fi\fi\fi

\ifdefined\nohyperref\else\ifdefined\hypersetup
  \definecolor{mydarkblue}{rgb}{0,0.08,0.45}
  \hypersetup{ %
    pdftitle={},
    pdfsubject={Proceedings of the International Conference on Machine Learning 2026},
    pdfkeywords={},
    pdfborder=0 0 0,
    pdfpagemode=UseNone,
    colorlinks=true,
    linkcolor=mydarkblue,
    citecolor=mydarkblue,
    filecolor=mydarkblue,
    urlcolor=mydarkblue,
    }
  \fi
\fi



% Uncomment the following for debugging.  It will cause LaTeX to dump
% the version of the ``appearing in'' string that will actually appear
% in the document.
%\typeout{>> Notice string='\Notice@String'}

% Change citation commands to be more like old ICML styles
\newcommand{\yrcite}[1]{\citeyearpar{#1}}
\renewcommand{\cite}[1]{\citep{#1}}


%%%%%%%%%%%%%%%%%%%%%%%%%%%%%%%%%%%%%%%%%%%%%%%%%%%%%%%%%%%%%%%%%%%%%%%%
% to ensure the letter format is used. pdflatex does compile the
% page size into the pdf. This is done using \pdfpagewidth and
% \pdfpageheight. As Latex does not know this directives, we first
% check whether pdflatex or latex is used.
%
% Kristian Kersting 2005
%
% in order to account for the more recent use of pdfetex as the default
% compiler, I have changed the pdf verification.
%
% Ricardo Silva 2007
%%%%%%%%%%%%%%%%%%%%%%%%%%%%%%%%%%%%%%%%%%%%%%%%%%%%%%%%%%%%%%%%%%%%%%%%%

\paperwidth=8.5in
\paperheight=11in

% old PDFLaTex verification, circa 2005
%
%\newif\ifpdf\ifx\pdfoutput\undefined
%  \pdffalse % we are not running PDFLaTeX
%\else
%  \pdfoutput=1 % we are running PDFLaTeX
%  \pdftrue
%\fi

\newif\ifpdf %adapted from ifpdf.sty
\ifx\pdfoutput\undefined
\else
   \ifx\pdfoutput\relax
   \else
     \ifcase\pdfoutput
     \else
       \pdftrue
     \fi
   \fi
\fi

\ifpdf
%    \pdfpagewidth=\paperwidth
%    \pdfpageheight=\paperheight
  \setlength{\pdfpagewidth}{8.5in}
  \setlength{\pdfpageheight}{11in}
\fi

% Physical page layout

\evensidemargin -0.23in
\oddsidemargin -0.23in
\setlength\textheight{9.0in}
\setlength\textwidth{6.75in}
\setlength\columnsep{0.25in}
\setlength\headheight{10pt}
\setlength\headsep{10pt}
\addtolength{\topmargin}{-20pt}
\addtolength{\topmargin}{-0.29in}

% Historically many authors tried to include packages like geometry or fullpage,
% which change the page layout. It either makes the proceedings inconsistent, or
% wastes organizers' time chasing authors. So let's nip these problems in the
% bud here. -- Iain Murray 2018.
%\RequirePackage{printlen}
\AtBeginDocument{%
\newif\ifmarginsmessedwith
\marginsmessedwithfalse
\ifdim\oddsidemargin=-16.62178pt     \else oddsidemargin has been altered.\\ \marginsmessedwithtrue\fi
\ifdim\headheight=10.0pt             \else headheight has been altered.\\ \marginsmessedwithtrue\fi
\ifdim\textheight=650.43pt           \else textheight has been altered.\\ \marginsmessedwithtrue\fi
\ifdim\marginparsep=11.0pt           \else marginparsep has been altered.\\ \marginsmessedwithtrue\fi
\ifdim\footskip=25.0pt               \else footskip has been altered.\\ \marginsmessedwithtrue\fi
\ifdim\hoffset=0.0pt                 \else hoffset has been altered.\\ \marginsmessedwithtrue\fi
\ifdim\paperwidth=614.295pt          \else paperwidth has been altered.\\ \marginsmessedwithtrue\fi
\ifdim\topmargin=-24.95781pt         \else topmargin has been altered.\\ \marginsmessedwithtrue\fi
\ifdim\headsep=10.0pt                \else headsep has been altered.\\ \marginsmessedwithtrue\fi
\ifdim\textwidth=487.8225pt          \else textwidth has been altered.\\ \marginsmessedwithtrue\fi
\ifdim\marginparpush=5.0pt           \else marginparpush has been altered.\\ \marginsmessedwithtrue\fi
\ifdim\voffset=0.0pt                 \else voffset has been altered.\\ \marginsmessedwithtrue\fi
\ifdim\paperheight=794.96999pt       \else paperheight has been altered.\\ \marginsmessedwithtrue\fi
\ifmarginsmessedwith

\textbf{\large \em The page layout violates the ICML style.}

Please do not change the page layout, or include packages like geometry,
savetrees, or fullpage, which change it for you.

We're not able to reliably undo arbitrary changes to the style. Please remove
the offending package(s), or layout-changing commands and try again.

\fi}


%% The following is adapted from code in the acmconf.sty conference
%% style file.  The constants in it are somewhat magical, and appear
%% to work well with the two-column format on US letter paper that
%% ICML uses, but will break if you change that layout, or if you use
%% a longer block of text for the copyright notice string.  Fiddle with
%% them if necessary to get the block to fit/look right.
%%
%% -- Terran Lane, 2003
%%
%% The following comments are included verbatim from acmconf.sty:
%%
%%% This section (written by KBT) handles the 1" box in the lower left
%%% corner of the left column of the first page by creating a picture,
%%% and inserting the predefined string at the bottom (with a negative
%%% displacement to offset the space allocated for a non-existent
%%% caption).
%%%
\def\ftype@copyrightbox{8}
\def\@copyrightspace{
\@float{copyrightbox}[b]
\begin{center}
\setlength{\unitlength}{1pc}
\begin{picture}(20,1.5)
\put(0,2.5){\line(1,0){4.818}}
\put(0,0){\parbox[b]{19.75pc}{\small \Notice@String}}
\end{picture}
\end{center}
\end@float}

\setlength\footskip{25.0pt}
\flushbottom \twocolumn
\sloppy

% Clear out the addcontentsline command
\def\addcontentsline#1#2#3{}

%%%%%%%%%%%%%%%%%%%%%%%%%%%%%%%%%%%%%%%%%%%%%%%%%%%%%%%%%%%%%%%%%%%%
%%% commands for formatting paper title, author names, and addresses.

% box to check the size of the running head
\newbox\titrun

% general page style
\pagestyle{fancy}
\fancyhf{}
\fancyfoot[C]{\thepage}
% set the width of the head rule to 1 point
\renewcommand{\headrulewidth}{1pt}

% definition to set the head as running head in the preamble
\def\icmltitlerunning#1{\gdef\@icmltitlerunning{#1}}

% main definition adapting \icmltitle from 2004
\long\def\icmltitle#1{%

   %check whether @icmltitlerunning exists
   % if not \icmltitle is used as running head
   \ifx\undefined\@icmltitlerunning%
      \gdef\@icmltitlerunning{#1}
   \fi

   %add it to pdf information
  \ifdefined\nohyperref\else\ifdefined\hypersetup
     \hypersetup{pdftitle={#1}}
   \fi\fi

   %get the dimension of the running title
   \global\setbox\titrun=\vbox{\small\bf\@icmltitlerunning}

   % error flag
   \gdef\@runningtitleerror{0}

    % running title too long
    \ifdim\wd\titrun>\textwidth%
      \gdef\@runningtitleerror{1}%
      % running title breaks a line
    \else \ifdim\ht\titrun>6.25pt
    \gdef\@runningtitleerror{2}%
      \fi
    \fi

       % if there is somthing wrong with the running title
    \ifnum\@runningtitleerror>0
      \typeout{}%
                 \typeout{}%
                 \typeout{*******************************************************}%
                 \typeout{Title exceeds size limitations for running head.}%
                 \typeout{Please supply a shorter form for the running head}
                 \typeout{with \string\icmltitlerunning{...}\space prior to \string\begin{document}}%
      \typeout{*******************************************************}%
      \typeout{}%
      \typeout{}%
      % set default running title
      \gdef\@icmltitlerunning{Title Suppressed Due to Excessive Size}
    \fi

    % no running title on the first page of the paper
    \thispagestyle{plain}

    {\center\baselineskip 18pt
      \toptitlebar{\Large\bf #1}\bottomtitlebar}
}

% set running title header
\fancyhead[C]{\small\bf\@icmltitlerunning}

\gdef\icmlfullauthorlist{}
\newcommand\addstringtofullauthorlist{\g@addto@macro\icmlfullauthorlist}
\newcommand\addtofullauthorlist[1]{%
  \ifdefined\icmlanyauthors%
    \addstringtofullauthorlist{, #1}%
  \else%
    \addstringtofullauthorlist{#1}%
    \gdef\icmlanyauthors{1}%
  \fi%
  \ifdefined\hypersetup%
    \hypersetup{pdfauthor=\icmlfullauthorlist}%
  \fi
}

\def\toptitlebar{\hrule height1pt \vskip .25in}
\def\bottomtitlebar{\vskip .22in \hrule height1pt \vskip .3in}

\newenvironment{icmlauthorlist}{%
  \setlength\topsep{0pt}
  \setlength\parskip{0pt}
  \begin{center}
    }{%
  \end{center}
}

\newcounter{@affiliationcounter}
\newcommand{\@pa}[1]{%
  \ifcsname the@affil#1\endcsname
    % do nothing
  \else
    \ifcsname @icmlsymbol#1\endcsname
      % nothing
    \else
      \stepcounter{@affiliationcounter}%
      \newcounter{@affil#1}%
      \setcounter{@affil#1}{\value{@affiliationcounter}}%
    \fi
  \fi%
  \ifcsname @icmlsymbol#1\endcsname
    \textsuperscript{\csname @icmlsymbol#1\endcsname\,}%
  \else
    \textsuperscript{\arabic{@affil#1}\,}%
  \fi
}

\newcommand{\icmlauthor}[2]{%
  \ificmlshowauthors
    \mbox{\bf #1}\,\@for\theaffil:=#2\do{\@pa{\theaffil}} \addtofullauthorlist{#1}%
  \else
    \ifdefined\@icmlfirsttime\else
      \gdef\@icmlfirsttime{1}
      \mbox{\bf Anonymous Authors}\@pa{@anon} \addtofullauthorlist{Anonymous Authors}
    \fi
  \fi
}

\newcommand{\icmlsetsymbol}[2]{%
  \expandafter\gdef\csname @icmlsymbol#1\endcsname{#2}
}

\newcommand{\icmlaffiliation}[2]{%
  \ificmlshowauthors
    \ifcsname the@affil#1\endcsname
      \expandafter\gdef\csname @affilname\csname the@affil#1\endcsname\endcsname{#2}%
    \else
      {\bf AUTHORERR: Error in use of \textbackslash{}icmlaffiliation command. Label ``#1'' not mentioned in some \textbackslash{}icmlauthor\{author name\}\{labels here\} command beforehand. }
      \typeout{}%
      \typeout{}%
      \typeout{*******************************************************}%
      \typeout{Affiliation label undefined. }%
      \typeout{Make sure \string\icmlaffiliation\space follows }%
      \typeout{all of \string\icmlauthor\space commands}%
      \typeout{*******************************************************}%
      \typeout{}%
      \typeout{}%
    \fi
  \else
    \expandafter\gdef\csname @affilname1\endcsname{Anonymous Institution, Anonymous City, Anonymous Region, Anonymous Country}
  \fi
}

\newcommand{\icmlcorrespondingauthor}[2]{%
  \ificmlshowauthors
    \ifdefined\icmlcorrespondingauthor@text
      \g@addto@macro\icmlcorrespondingauthor@text{, #1 \textless{}#2\textgreater{}}
    \else
      \gdef\icmlcorrespondingauthor@text{#1 \textless{}#2\textgreater{}}
    \fi
  \else
    \gdef\icmlcorrespondingauthor@text{Anonymous Author \textless{}anon.email@domain.com\textgreater{}}
  \fi
}

\newcommand{\icmlEqualContribution}{\textsuperscript{*}Equal contribution }


% --- ICML 2026: ensure authors do not omit the affiliations/notice footnote ---
\newif\ificml@noticeprinted
\icml@noticeprintedfalse
\AtEndDocument{%
  \ificml@noticeprinted\relax\else
    \PackageWarningNoLine{icml2026}{%
      You did not call \string\printAffiliationsAndNotice{}. If you have no notice,%
      call \string\printAffiliationsAndNotice\string{} (empty braces).%
    }%
  \fi
}


\newcounter{@affilnum}
\newcommand{\printAffiliationsAndNotice}[1]{\global\icml@noticeprintedtrue%
  \stepcounter{@affiliationcounter}%
  {\let\thefootnote\relax\footnotetext{\hspace*{-\footnotesep}\ificmlshowauthors #1\fi%
      \forloop{@affilnum}{1}{\value{@affilnum} < \value{@affiliationcounter}}{
        \textsuperscript{\arabic{@affilnum}}\ifcsname @affilname\the@affilnum\endcsname%
          \csname @affilname\the@affilnum\endcsname%
        \else
          {\bf AUTHORERR: Missing \textbackslash{}icmlaffiliation.}
        \fi
      }.%
      \ifdefined\icmlcorrespondingauthor@text
         { }Correspondence to: \icmlcorrespondingauthor@text.
      \else
        {\bf AUTHORERR: Missing \textbackslash{}icmlcorrespondingauthor.}
      \fi

      \ \\
      \Notice@String
    }
  }
}

\long\def\icmladdress#1{%
  {\bf The \textbackslash{}icmladdress command is no longer used.  See the example\_paper PDF .tex for usage of \textbackslash{}icmlauther and \textbackslash{}icmlaffiliation.}
}

%% keywords as first class citizens
\def\icmlkeywords#1{%
  \ifdefined\nohyperref\else\ifdefined\hypersetup
      \hypersetup{pdfkeywords={#1}}
    \fi\fi
}

% modification to natbib citations
\setcitestyle{authoryear,round,citesep={;},aysep={,},yysep={;}}

% Redefinition of the abstract environment.
\renewenvironment{abstract}
{%
  \centerline{\large\bf Abstract}
  \vspace{-0.12in}\begin{quote}}
    {\par\end{quote}\vskip 0.12in}

% numbered section headings with different treatment of numbers

\def\@startsection#1#2#3#4#5#6{\if@noskipsec \leavevmode \fi
  \par \@tempskipa #4\relax
  \@afterindenttrue
  \ifdim \@tempskipa <\z@ \@tempskipa -\@tempskipa \fi
  \if@nobreak \everypar{}\else
    \addpenalty{\@secpenalty}\addvspace{\@tempskipa}\fi \@ifstar
  {\@ssect{#3}{#4}{#5}{#6}}{\@dblarg{\@sict{#1}{#2}{#3}{#4}{#5}{#6}}}}

\def\@sict#1#2#3#4#5#6[#7]#8{\ifnum #2>\c@secnumdepth
    \def\@svsec{}\else
    \refstepcounter{#1}\edef\@svsec{\csname the#1\endcsname}\fi
  \@tempskipa #5\relax
  \ifdim \@tempskipa>\z@
    \begingroup #6\relax
    \@hangfrom{\hskip #3\relax\@svsec.~}{\interlinepenalty \@M #8\par}
    \endgroup
    \csname #1mark\endcsname{#7}\addcontentsline
    {toc}{#1}{\ifnum #2>\c@secnumdepth \else
        \protect\numberline{\csname the#1\endcsname}\fi
      #7}\else
    \def\@svsechd{#6\hskip #3\@svsec #8\csname #1mark\endcsname
      {#7}\addcontentsline
      {toc}{#1}{\ifnum #2>\c@secnumdepth \else
          \protect\numberline{\csname the#1\endcsname}\fi
        #7}}\fi
  \@xsect{#5}}

\def\@sect#1#2#3#4#5#6[#7]#8{\ifnum #2>\c@secnumdepth
    \def\@svsec{}\else
    \refstepcounter{#1}\edef\@svsec{\csname the#1\endcsname\hskip 0.4em }\fi
  \@tempskipa #5\relax
  \ifdim \@tempskipa>\z@
    \begingroup #6\relax
    \@hangfrom{\hskip #3\relax\@svsec}{\interlinepenalty \@M #8\par}
    \endgroup
    \csname #1mark\endcsname{#7}\addcontentsline
    {toc}{#1}{\ifnum #2>\c@secnumdepth \else
        \protect\numberline{\csname the#1\endcsname}\fi
      #7}\else
    \def\@svsechd{#6\hskip #3\@svsec #8\csname #1mark\endcsname
      {#7}\addcontentsline
      {toc}{#1}{\ifnum #2>\c@secnumdepth \else
          \protect\numberline{\csname the#1\endcsname}\fi
        #7}}\fi
  \@xsect{#5}}

% section headings with less space above and below them
\def\thesection {\arabic{section}}
\def\thesubsection {\thesection.\arabic{subsection}}
\def\section{\@startsection{section}{1}{\z@}{-0.12in}{0.02in}
  {\large\bf\raggedright}}
\def\subsection{\@startsection{subsection}{2}{\z@}{-0.10in}{0.01in}
  {\normalsize\bf\raggedright}}
\def\subsubsection{\@startsection{subsubsection}{3}{\z@}{-0.08in}{0.01in}
  {\normalsize\sc\raggedright}}
\def\paragraph{\@startsection{paragraph}{4}{\z@}{1.5ex plus
    0.5ex minus .2ex}{-1em}{\normalsize\bf}}
\def\subparagraph{\@startsection{subparagraph}{5}{\z@}{1.5ex plus
    0.5ex minus .2ex}{-1em}{\normalsize\bf}}

% Footnotes
\footnotesep 6.65pt %
\skip\footins 9pt
\def\footnoterule{\kern-3pt \hrule width 0.8in \kern 2.6pt }
\setcounter{footnote}{0}

% Lists and paragraphs
\parindent 0pt
\topsep 4pt plus 1pt minus 2pt
\partopsep 1pt plus 0.5pt minus 0.5pt
\itemsep 2pt plus 1pt minus 0.5pt
\parsep 2pt plus 1pt minus 0.5pt
\parskip 6pt

\leftmargin 2em \leftmargini\leftmargin \leftmarginii 2em
\leftmarginiii 1.5em \leftmarginiv 1.0em \leftmarginv .5em
\leftmarginvi .5em
\labelwidth\leftmargini\advance\labelwidth-\labelsep \labelsep 5pt

\def\@listi{\leftmargin\leftmargini}
\def\@listii{\leftmargin\leftmarginii
  \labelwidth\leftmarginii\advance\labelwidth-\labelsep
  \topsep 2pt plus 1pt minus 0.5pt
  \parsep 1pt plus 0.5pt minus 0.5pt
  \itemsep \parsep}
\def\@listiii{\leftmargin\leftmarginiii
  \labelwidth\leftmarginiii\advance\labelwidth-\labelsep
  \topsep 1pt plus 0.5pt minus 0.5pt
  \parsep \z@ \partopsep 0.5pt plus 0pt minus 0.5pt
  \itemsep \topsep}
\def\@listiv{\leftmargin\leftmarginiv
  \labelwidth\leftmarginiv\advance\labelwidth-\labelsep}
\def\@listv{\leftmargin\leftmarginv
  \labelwidth\leftmarginv\advance\labelwidth-\labelsep}
\def\@listvi{\leftmargin\leftmarginvi
  \labelwidth\leftmarginvi\advance\labelwidth-\labelsep}

\abovedisplayskip 7pt plus2pt minus5pt%
\belowdisplayskip \abovedisplayskip
\abovedisplayshortskip  0pt plus3pt%
\belowdisplayshortskip  4pt plus3pt minus3pt%

% Less leading in most fonts (due to the narrow columns)
% The choices were between 1-pt and 1.5-pt leading
\def\@normalsize{\@setsize\normalsize{11pt}\xpt\@xpt}
\def\small{\@setsize\small{10pt}\ixpt\@ixpt}
\def\footnotesize{\@setsize\footnotesize{10pt}\ixpt\@ixpt}
\def\scriptsize{\@setsize\scriptsize{8pt}\viipt\@viipt}
\def\tiny{\@setsize\tiny{7pt}\vipt\@vipt}
\def\large{\@setsize\large{14pt}\xiipt\@xiipt}
\def\Large{\@setsize\Large{16pt}\xivpt\@xivpt}
\def\LARGE{\@setsize\LARGE{20pt}\xviipt\@xviipt}
\def\huge{\@setsize\huge{23pt}\xxpt\@xxpt}
\def\Huge{\@setsize\Huge{28pt}\xxvpt\@xxvpt}

% Revised formatting for figure captions and table titles.
\captionsetup{
  skip=0.1in,
  font=small,
  labelfont={it,small},
  labelsep=period
}
\captionsetup[table]{position=above}
\captionsetup[figure]{position=below}

\def\fnum@figure{Figure \thefigure}
\def\fnum@table{Table \thetable}

% Strut macros for skipping spaces above and below text in tables.
\def\abovestrut#1{\rule[0in]{0in}{#1}\ignorespaces}
\def\belowstrut#1{\rule[-#1]{0in}{#1}\ignorespaces}

\def\abovespace{\abovestrut{0.20in}}
\def\aroundspace{\abovestrut{0.20in}\belowstrut{0.10in}}
\def\belowspace{\belowstrut{0.10in}}

% Various personal itemization commands.
\def\texitem#1{\par\noindent\hangindent 12pt
  \hbox to 12pt {\hss #1 ~}\ignorespaces}
\def\icmlitem{\texitem{$\bullet$}}

% To comment out multiple lines of text.
\long\def\comment#1{}

%% Line counter (not in final version). Adapted from NIPS style file by Christoph Sawade

% Vertical Ruler
% This code is, largely, from the CVPR 2010 conference style file
% ----- define vruler
\makeatletter
\newbox\icmlrulerbox
\newcount\icmlrulercount
\newdimen\icmlruleroffset
\newdimen\cv@lineheight
\newdimen\cv@boxheight
\newbox\cv@tmpbox
\newcount\cv@refno
\newcount\cv@tot
% NUMBER with left flushed zeros  \fillzeros[<WIDTH>]<NUMBER>
\newcount\cv@tmpc@ \newcount\cv@tmpc
\def\fillzeros[#1]#2{\cv@tmpc@=#2\relax\ifnum\cv@tmpc@<0\cv@tmpc@=-\cv@tmpc@\fi
  \cv@tmpc=1 %
  \loop\ifnum\cv@tmpc@<10 \else \divide\cv@tmpc@ by 10 \advance\cv@tmpc by 1 \fi
  \ifnum\cv@tmpc@=10\relax\cv@tmpc@=11\relax\fi \ifnum\cv@tmpc@>10 \repeat
  \ifnum#2<0\advance\cv@tmpc1\relax-\fi
  \loop\ifnum\cv@tmpc<#1\relax0\advance\cv@tmpc1\relax\fi \ifnum\cv@tmpc<#1 \repeat
  \cv@tmpc@=#2\relax\ifnum\cv@tmpc@<0\cv@tmpc@=-\cv@tmpc@\fi \relax\the\cv@tmpc@}%
% \makevruler[<SCALE>][<INITIAL_COUNT>][<STEP>][<DIGITS>][<HEIGHT>]
\def\makevruler[#1][#2][#3][#4][#5]{
  \begingroup\offinterlineskip
  \textheight=#5\vbadness=10000\vfuzz=120ex\overfullrule=0pt%
  \global\setbox\icmlrulerbox=\vbox to \textheight{%
    {
        \parskip=0pt\hfuzz=150em\cv@boxheight=\textheight
        \cv@lineheight=#1\global\icmlrulercount=#2%
        \cv@tot\cv@boxheight\divide\cv@tot\cv@lineheight\advance\cv@tot2%
        \cv@refno1\vskip-\cv@lineheight\vskip1ex%
        \loop\setbox\cv@tmpbox=\hbox to0cm{\hfil {\hfil\fillzeros[#4]\icmlrulercount}}%
        \ht\cv@tmpbox\cv@lineheight\dp\cv@tmpbox0pt\box\cv@tmpbox\break
        \advance\cv@refno1\global\advance\icmlrulercount#3\relax
        \ifnum\cv@refno<\cv@tot\repeat
      }
  }
  \endgroup
}%
\makeatother
% ----- end of vruler

% \makevruler[<SCALE>][<INITIAL_COUNT>][<STEP>][<DIGITS>][<HEIGHT>]
\def\icmlruler#1{\makevruler[12pt][#1][1][3][\textheight]\usebox{\icmlrulerbox}}
\AddToShipoutPicture{%
  \icmlruleroffset=\textheight
  \advance\icmlruleroffset by 5.2pt % top margin
  \color[rgb]{.7,.7,.7}
  \ificmlshowauthors\else
    \AtTextUpperLeft{%
      \put(\LenToUnit{-35pt},\LenToUnit{-\icmlruleroffset}){%left ruler
        \icmlruler{\icmlrulercount}}
      %\put(\LenToUnit{1.04\textwidth},\LenToUnit{-\icmlruleroffset}){%right ruler
      %  \icmlruler{\icmlrulercount}}
    }
  \fi
}
\endinput






\end{filecontents*}
\documentclass{article}
\usepackage{luatex85}
\usepackage[T1]{fontenc}
\usepackage{microtype,graphicx,booktabs}
\usepackage{hyperref,xurl}
\usepackage[accepted]{icml2026}
% Author-visible ICML typography, without a preprint or proceedings notice.
\makeatletter\renewcommand{\Notice@String}{}\makeatother
\usepackage{amsmath,amssymb,tikz}
\definecolor{panelnavy}{HTML}{081A32}
\definecolor{treeink}{HTML}{EDF5FF}
\definecolor{rootpurple}{HTML}{FFFFFF}
\definecolor{leafteal}{HTML}{EDF5FF}
\definecolor{leafgray}{HTML}{EDF5FF}
\definecolor{treeblue}{HTML}{00E5FF}
\definecolor{treeorange}{HTML}{FFEE32}
\hypersetup{pdftitle={Compact Decision Programs over Frozen Neural Features for Lesion Segmentation},pdfauthor={Brian Mingus},hidelinks}
\icmltitlerunning{Compact Decision Programs over Frozen Neural Features}
\newcommand{\doi}[1]{\href{https://doi.org/#1}{doi:~\nolinkurl{#1}}}
\begin{document}
\twocolumn[
\icmltitle{Compact Decision Programs over Frozen Neural Features\\for Lesion Segmentation}
\begin{icmlauthorlist}
\icmlauthor{Brian Mingus}{cognatory}
\end{icmlauthorlist}
\icmlaffiliation{cognatory}{Cognatory, LLC, Denver, Colorado, USA}
\icmlcorrespondingauthor{Brian Mingus}{brian@cognatory.com}
\icmlkeywords{Decision programs, gradient boosting, segmentation, frozen features}
\vskip 0.12in
{\begin{minipage}{.94\textwidth}\small Frozen neural features support compact decision programs for location-assisted lesion segmentation. On 32 RIDER patients and 62 scans, patient-grouped evaluation yields mean Dice of \textbf{0.765}, versus \textbf{0.734} for the pretrained mask and \textbf{0.758} for an earlier hybrid procedure. Current and earlier all-patient refits yield complete binary trees with 181 and 21 nodes, respectively. Both trees are shown in full. The neural feature pipeline remains necessary; these exploratory results demonstrate adaptable, inspectable downstream decisions.\end{minipage}}
\vskip 0.12in
{\begin{minipage}{\textwidth}
\centering
{\setlength{\fboxsep}{7pt}\colorbox{panelnavy}{\begin{minipage}{\dimexpr\textwidth-14pt\relax}\centering\color{white}
\begin{minipage}[t]{.30\textwidth}\centering\vspace{0pt}
{\small\bfseries (a) Earlier: Dice 0.758}\\{\footnotesize 21 nodes; $\sim$250 B}\par\smallskip
\begin{tikzpicture}[x=1in,y=1in, every node/.style={inner sep=.6pt,outer sep=0pt,font=\fontsize{7}{8}\selectfont},
trueedge/.style={draw=treeblue,line width=.6pt,preaction={draw=treeblue,opacity=.08,line width=1.4pt}},falseedge/.style={draw=treeorange,densely dashed,line width=.6pt,preaction={draw=treeorange,opacity=.08,line width=1.4pt}}]
\node[text=leafgray] (n1) at (0.61458,1.94970) {$\mathbf{B}$};
\node[text=leafteal] (n5) at (0.43403,1.14760) {$\mathbf{L}$};
\node[text=leafgray] (n10) at (0.10069,0.07812) {$\mathbf{B}$};
\node[text=leafteal] (n11) at (0.22569,0.07812) {$\mathbf{L}$};
\node[text=treeink] (n9) at (0.15625,0.34549) {$x<\tau_{7}$};
\node[text=leafteal] (n12) at (0.39236,0.34549) {$\mathbf{L}$};
\node[text=treeink] (n8) at (0.32292,0.61285) {$u<\tau_{6}$};
\node[text=leafgray] (n14) at (0.62847,0.34549) {$\mathbf{B}$};
\node[text=leafgray] (n16) at (0.80903,0.07812) {$\mathbf{B}$};
\node[text=leafteal] (n17) at (0.93403,0.07812) {$\mathbf{L}$};
\node[text=treeink] (n15) at (0.87847,0.34549) {$x<\tau_{7}$};
\node[text=treeink] (n13) at (0.68403,0.61285) {$u<\tau_{6}$};
\node[text=treeink] (n7) at (0.50347,0.88021) {$a_{21}<\tau_{5}$};
\node[text=leafteal] (n18) at (0.86458,0.88021) {$\mathbf{L}$};
\node[text=treeink] (n6) at (0.68403,1.14760) {$u<\tau_{4}$};
\node[text=treeink] (n4) at (0.55903,1.41490) {$a_{21}<\tau_{3}$};
\node[text=leafteal] (n19) at (0.92014,1.41490) {$\mathbf{L}$};
\node[text=treeink] (n3) at (0.73958,1.68230) {$x<\tau_{2}$};
\node[text=leafteal] (n20) at (0.97569,1.68230) {$\mathbf{L}$};
\node[text=treeink] (n2) at (0.86458,1.94970) {$x<\tau_{1}$};
\node[text=rootpurple] (n0) at (0.73958,2.21700) {$m<1$};
\draw[trueedge] (n0.south) -- (n1.north);
\draw[falseedge] (n0.south) -- (n2.north);
\draw[trueedge] (n2.south) -- (n3.north);
\draw[trueedge] (n3.south) -- (n4.north);
\draw[trueedge] (n4.south) -- (n5.north);
\draw[falseedge] (n4.south) -- (n6.north);
\draw[trueedge] (n6.south) -- (n7.north);
\draw[trueedge] (n7.south) -- (n8.north);
\draw[trueedge] (n8.south) -- (n9.north);
\draw[trueedge] (n9.south) -- (n10.north);
\draw[falseedge] (n9.south) -- (n11.north);
\draw[falseedge] (n8.south) -- (n12.north);
\draw[falseedge] (n7.south) -- (n13.north);
\draw[trueedge] (n13.south) -- (n14.north);
\draw[falseedge] (n13.south) -- (n15.north);
\draw[trueedge] (n15.south) -- (n16.north);
\draw[falseedge] (n15.south) -- (n17.north);
\draw[falseedge] (n6.south) -- (n18.north);
\draw[falseedge] (n3.south) -- (n19.north);
\draw[falseedge] (n2.south) -- (n20.north);
\end{tikzpicture}
\end{minipage}\hfill\begin{minipage}[t]{.68\textwidth}\centering\vspace{0pt}
{\small\bfseries (b) Current: Dice 0.765}\\{\footnotesize 181 nodes; $\sim$770 B}\par\smallskip
\begin{tikzpicture}[x=1in,y=1in, every node/.style={inner sep=.6pt,outer sep=0pt,font=\fontsize{6.5}{7.5}\selectfont},
trueedge/.style={draw=treeblue,line width=.6pt,preaction={draw=treeblue,opacity=.08,line width=1.4pt}},falseedge/.style={draw=treeorange,densely dashed,line width=.6pt,preaction={draw=treeorange,opacity=.08,line width=1.4pt}}]
\node[text=leafgray] (n2) at (1.58728,1.87820) {$\mathbf{B}$};
\node[text=leafgray] (n4) at (1.58728,1.73668) {$\mathbf{B}$};
\node[text=leafteal] (n6) at (1.61523,1.59521) {$\mathbf{L}$};
\node[text=leafgray] (n7) at (1.73056,1.59521) {$\mathbf{B}$};
\node[text=treeink] (n5) at (1.74439,1.73668) {$P_{4}$};
\node[text=treeink] (n3) at (1.74439,1.87820) {$P_{3}$};
\node[text=treeink] (n1) at (1.74439,2.01967) {$P_{2}$};
\node[text=leafteal] (n12) at (1.74439,1.45373) {$\mathbf{L}$};
\node[text=leafteal] (n17) at (0.48338,0.88779) {$\mathbf{L}$};
\node[text=leafgray] (n23) at (0.03152,0.18035) {$\mathbf{B}$};
\node[text=leafteal] (n24) at (0.12610,0.18035) {$\mathbf{L}$};
\node[text=treeink] (n22) at (0.10508,0.32183) {$P_{17}$};
\node[text=leafgray] (n27) at (0.24169,0.03886) {$\mathbf{B}$};
\node[text=leafteal] (n28) at (0.33626,0.03886) {$\mathbf{L}$};
\node[text=treeink] (n26) at (0.29423,0.18035) {$P_{17}$};
\node[text=leafteal] (n29) at (0.44134,0.18035) {$\mathbf{L}$};
\node[text=treeink] (n25) at (0.30474,0.32183) {$P_{18}$};
\node[text=treeink] (n21) at (0.19966,0.46332) {$P_{16}$};
\node[text=leafteal] (n30) at (0.34677,0.46332) {$\mathbf{L}$};
\node[text=treeink] (n20) at (0.31525,0.60479) {$P_{15}$};
\node[text=leafgray] (n32) at (0.49389,0.46332) {$\mathbf{B}$};
\node[text=leafteal] (n33) at (0.58846,0.46332) {$\mathbf{L}$};
\node[text=treeink] (n31) at (0.51491,0.60479) {$P_{18}$};
\node[text=treeink] (n19) at (0.48338,0.74627) {$P_{14}$};
\node[text=leafteal] (n34) at (0.63049,0.74627) {$\mathbf{L}$};
\node[text=treeink] (n18) at (0.63049,0.88779) {$P_{13}$};
\node[text=treeink] (n16) at (0.63049,1.02926) {$P_{12}$};
\node[text=leafteal] (n35) at (0.83013,1.02926) {$\mathbf{L}$};
\node[text=treeink] (n15) at (0.83013,1.17074) {$P_{11}$};
\node[text=leafteal] (n36) at (1.07186,1.17074) {$\mathbf{L}$};
\node[text=treeink] (n14) at (1.07186,1.31226) {$P_{10}$};
\node[text=leafteal] (n37) at (1.85993,1.31226) {$\mathbf{L}$};
\node[text=treeink] (n13) at (1.85993,1.45373) {$P_{9}$};
\node[text=treeink] (n11) at (1.85993,1.59521) {$P_{8}$};
\node[text=leafteal] (n38) at (1.97553,1.59521) {$\mathbf{L}$};
\node[text=treeink] (n10) at (1.92303,1.73668) {$P_{7}$};
\node[text=leafgray] (n47) at (0.71455,0.60479) {$\mathbf{B}$};
\node[text=leafteal] (n48) at (0.80910,0.60479) {$\mathbf{L}$};
\node[text=treeink] (n46) at (0.77762,0.74627) {$P_{17}$};
\node[text=leafteal] (n49) at (0.92470,0.74627) {$\mathbf{L}$};
\node[text=treeink] (n45) at (0.92470,0.88779) {$P_{15}$};
\node[text=leafteal] (n50) at (1.07186,0.88779) {$\mathbf{L}$};
\node[text=treeink] (n44) at (1.07186,1.02926) {$P_{13}$};
\node[text=leafgray] (n54) at (0.92470,0.60479) {$\mathbf{B}$};
\node[text=leafgray] (n57) at (0.75659,0.32183) {$\mathbf{B}$};
\node[text=leafgray] (n59) at (0.86168,0.18035) {$\mathbf{B}$};
\node[text=leafteal] (n60) at (0.95625,0.18035) {$\mathbf{L}$};
\node[text=treeink] (n58) at (0.90367,0.32183) {$P_{18}$};
\node[text=treeink] (n56) at (0.88271,0.46332) {$P_{17}$};
\node[text=leafgray] (n62) at (1.06134,0.32183) {$\mathbf{B}$};
\node[text=leafgray] (n64) at (1.15592,0.18035) {$\mathbf{B}$};
\node[text=leafteal] (n65) at (1.25049,0.18035) {$\mathbf{L}$};
\node[text=treeink] (n63) at (1.20843,0.32183) {$P_{17}$};
\node[text=treeink] (n61) at (1.08238,0.46332) {$P_{18}$};
\node[text=treeink] (n55) at (1.07186,0.60479) {$P_{16}$};
\node[text=treeink] (n53) at (1.07186,0.74627) {$P_{15}$};
\node[text=leafgray] (n66) at (1.21894,0.74627) {$\mathbf{B}$};
\node[text=treeink] (n52) at (1.21894,0.88779) {$P_{14}$};
\node[text=leafgray] (n69) at (1.21894,0.60479) {$\mathbf{B}$};
\node[text=leafteal] (n71) at (1.30300,0.46332) {$\mathbf{L}$};
\node[text=leafgray] (n73) at (1.40809,0.32183) {$\mathbf{B}$};
\node[text=leafteal] (n74) at (1.50266,0.32183) {$\mathbf{L}$};
\node[text=treeink] (n72) at (1.45016,0.46332) {$P_{18}$};
\node[text=treeink] (n70) at (1.36610,0.60479) {$P_{16}$};
\node[text=treeink] (n68) at (1.36610,0.74627) {$P_{15}$};
\node[text=leafgray] (n75) at (1.51318,0.74627) {$\mathbf{B}$};
\node[text=treeink] (n67) at (1.41861,0.88779) {$P_{14}$};
\node[text=treeink] (n51) at (1.34507,1.02926) {$P_{13}$};
\node[text=treeink] (n43) at (1.34507,1.17074) {$P_{12}$};
\node[text=leafgray] (n80) at (1.51318,0.60479) {$\mathbf{B}$};
\node[text=leafteal] (n81) at (1.60775,0.60479) {$\mathbf{L}$};
\node[text=treeink] (n79) at (1.66026,0.74627) {$P_{17}$};
\node[text=leafteal] (n82) at (1.80742,0.74627) {$\mathbf{L}$};
\node[text=treeink] (n78) at (1.80742,0.88779) {$P_{15}$};
\node[text=leafteal] (n83) at (1.95450,0.88779) {$\mathbf{L}$};
\node[text=treeink] (n77) at (1.95450,1.02926) {$P_{13}$};
\node[text=leafgray] (n88) at (1.62879,0.46332) {$\mathbf{B}$};
\node[text=leafgray] (n90) at (1.66026,0.32183) {$\mathbf{B}$};
\node[text=leafgray] (n92) at (1.73388,0.18035) {$\mathbf{B}$};
\node[text=leafteal] (n93) at (1.82845,0.18035) {$\mathbf{L}$};
\node[text=treeink] (n91) at (1.80742,0.32183) {$P_{17}$};
\node[text=treeink] (n89) at (1.77587,0.46332) {$P_{18}$};
\node[text=treeink] (n87) at (1.75484,0.60479) {$P_{16}$};
\node[text=leafteal] (n95) at (1.94398,0.46332) {$\mathbf{L}$};
\node[text=leafgray] (n98) at (1.98605,0.18035) {$\mathbf{B}$};
\node[text=leafteal] (n99) at (2.08062,0.18035) {$\mathbf{L}$};
\node[text=treeink] (n97) at (2.03856,0.32183) {$P_{17}$};
\node[text=leafteal] (n100) at (2.18571,0.32183) {$\mathbf{L}$};
\node[text=treeink] (n96) at (2.09114,0.46332) {$P_{18}$};
\node[text=treeink] (n94) at (1.95450,0.60479) {$P_{16}$};
\node[text=treeink] (n86) at (1.95450,0.74627) {$P_{15}$};
\node[text=leafgray] (n102) at (2.10166,0.60479) {$\mathbf{B}$};
\node[text=leafgray] (n104) at (2.24874,0.46332) {$\mathbf{B}$};
\node[text=leafteal] (n105) at (2.34331,0.46332) {$\mathbf{L}$};
\node[text=treeink] (n103) at (2.24874,0.60479) {$P_{17}$};
\node[text=treeink] (n101) at (2.15416,0.74627) {$P_{18}$};
\node[text=treeink] (n85) at (2.15416,0.88779) {$P_{14}$};
\node[text=leafgray] (n109) at (2.43788,0.46332) {$\mathbf{B}$};
\node[text=leafteal] (n110) at (2.53246,0.46332) {$\mathbf{L}$};
\node[text=treeink] (n108) at (2.44840,0.60479) {$P_{18}$};
\node[text=leafteal] (n111) at (2.59556,0.60479) {$\mathbf{L}$};
\node[text=treeink] (n107) at (2.46943,0.74627) {$P_{15}$};
\node[text=leafgray] (n113) at (2.69013,0.60479) {$\mathbf{B}$};
\node[text=leafteal] (n114) at (2.78471,0.60479) {$\mathbf{L}$};
\node[text=treeink] (n112) at (2.69013,0.74627) {$P_{18}$};
\node[text=treeink] (n106) at (2.46943,0.88779) {$P_{14}$};
\node[text=treeink] (n84) at (2.15416,1.02926) {$P_{13}$};
\node[text=treeink] (n76) at (2.05959,1.17074) {$P_{12}$};
\node[text=treeink] (n42) at (2.05959,1.31226) {$P_{11}$};
\node[text=leafteal] (n117) at (2.77419,1.02926) {$\mathbf{L}$};
\node[text=leafteal] (n119) at (2.84773,0.88779) {$\mathbf{L}$};
\node[text=leafgray] (n124) at (2.80566,0.32183) {$\mathbf{B}$};
\node[text=leafteal] (n125) at (2.90024,0.32183) {$\mathbf{L}$};
\node[text=treeink] (n123) at (2.85825,0.46332) {$P_{11}$};
\node[text=leafteal] (n126) at (3.00533,0.46332) {$\mathbf{L}$};
\node[text=treeink] (n122) at (2.93179,0.60479) {$P_{17}$};
\node[text=leafteal] (n127) at (3.07887,0.60479) {$\mathbf{L}$};
\node[text=treeink] (n121) at (2.93179,0.74627) {$P_{20}$};
\node[text=leafteal] (n128) at (3.07887,0.74627) {$\mathbf{L}$};
\node[text=treeink] (n120) at (2.99481,0.88779) {$P_{15}$};
\node[text=treeink] (n118) at (2.92127,1.02926) {$P_{19}$};
\node[text=treeink] (n116) at (2.84773,1.17074) {$P_{12}$};
\node[text=leafteal] (n131) at (3.18396,0.88779) {$\mathbf{L}$};
\node[text=leafgray] (n135) at (3.15248,0.46332) {$\mathbf{B}$};
\node[text=leafteal] (n136) at (3.24706,0.46332) {$\mathbf{L}$};
\node[text=treeink] (n134) at (3.22602,0.60479) {$P_{11}$};
\node[text=leafteal] (n137) at (3.37311,0.60479) {$\mathbf{L}$};
\node[text=treeink] (n133) at (3.22602,0.74627) {$P_{17}$};
\node[text=leafteal] (n138) at (3.37311,0.74627) {$\mathbf{L}$};
\node[text=treeink] (n132) at (3.33112,0.88779) {$P_{20}$};
\node[text=treeink] (n130) at (3.26802,1.02926) {$P_{14}$};
\node[text=leafgray] (n144) at (3.41517,0.32183) {$\mathbf{B}$};
\node[text=leafteal] (n145) at (3.50975,0.32183) {$\mathbf{L}$};
\node[text=treeink] (n143) at (3.45716,0.46332) {$P_{11}$};
\node[text=leafteal] (n146) at (3.60432,0.46332) {$\mathbf{L}$};
\node[text=treeink] (n142) at (3.52026,0.60479) {$P_{17}$};
\node[text=leafteal] (n147) at (3.66734,0.60479) {$\mathbf{L}$};
\node[text=treeink] (n141) at (3.52026,0.74627) {$P_{20}$};
\node[text=leafteal] (n148) at (3.66734,0.74627) {$\mathbf{L}$};
\node[text=treeink] (n140) at (3.55174,0.88779) {$P_{15}$};
\node[text=leafgray] (n152) at (3.75140,0.46332) {$\mathbf{B}$};
\node[text=leafteal] (n153) at (3.84598,0.46332) {$\mathbf{L}$};
\node[text=treeink] (n151) at (3.81450,0.60479) {$P_{11}$};
\node[text=leafteal] (n154) at (3.96158,0.60479) {$\mathbf{L}$};
\node[text=treeink] (n150) at (3.81450,0.74627) {$P_{17}$};
\node[text=leafteal] (n155) at (3.96158,0.74627) {$\mathbf{L}$};
\node[text=treeink] (n149) at (3.81450,0.88779) {$P_{20}$};
\node[text=treeink] (n139) at (3.55174,1.02926) {$P_{14}$};
\node[text=treeink] (n129) at (3.26802,1.17074) {$P_{19}$};
\node[text=treeink] (n115) at (2.84773,1.31226) {$P_{16}$};
\node[text=treeink] (n41) at (2.10166,1.45373) {$P_{10}$};
\node[text=leafteal] (n159) at (3.84598,1.02926) {$\mathbf{L}$};
\node[text=leafgray] (n166) at (3.87752,0.18035) {$\mathbf{B}$};
\node[text=leafteal] (n167) at (3.97210,0.18035) {$\mathbf{L}$};
\node[text=treeink] (n165) at (3.93003,0.32183) {$P_{18}$};
\node[text=leafteal] (n168) at (4.07719,0.32183) {$\mathbf{L}$};
\node[text=treeink] (n164) at (4.06667,0.46332) {$P_{17}$};
\node[text=leafgray] (n171) at (4.19280,0.18035) {$\mathbf{B}$};
\node[text=leafteal] (n172) at (4.28737,0.18035) {$\mathbf{L}$};
\node[text=treeink] (n170) at (4.24530,0.32183) {$P_{17}$};
\node[text=leafteal] (n173) at (4.39246,0.32183) {$\mathbf{L}$};
\node[text=treeink] (n169) at (4.26634,0.46332) {$P_{18}$};
\node[text=treeink] (n163) at (4.10874,0.60479) {$P_{16}$};
\node[text=leafteal] (n174) at (4.25582,0.60479) {$\mathbf{L}$};
\node[text=treeink] (n162) at (4.10874,0.74627) {$P_{15}$};
\node[text=leafteal] (n175) at (4.25582,0.74627) {$\mathbf{L}$};
\node[text=treeink] (n161) at (4.10874,0.88779) {$P_{14}$};
\node[text=leafteal] (n176) at (4.25582,0.88779) {$\mathbf{L}$};
\node[text=treeink] (n160) at (4.10874,1.02926) {$P_{13}$};
\node[text=treeink] (n158) at (3.84598,1.17074) {$P_{12}$};
\node[text=leafteal] (n177) at (3.99313,1.17074) {$\mathbf{L}$};
\node[text=treeink] (n157) at (3.84598,1.31226) {$P_{10}$};
\node[text=leafteal] (n178) at (3.99313,1.31226) {$\mathbf{L}$};
\node[text=treeink] (n156) at (3.84598,1.45373) {$P_{21}$};
\node[text=treeink] (n40) at (2.10166,1.59521) {$P_{9}$};
\node[text=leafteal] (n179) at (2.21726,1.59521) {$\mathbf{L}$};
\node[text=treeink] (n39) at (2.10166,1.73668) {$P_{7}$};
\node[text=treeink] (n9) at (1.92303,1.87820) {$P_{6}$};
\node[text=leafteal] (n180) at (2.03856,1.87820) {$\mathbf{L}$};
\node[text=treeink] (n8) at (1.92303,2.01967) {$P_{5}$};
\node[text=rootpurple] (n0) at (1.83897,2.16115) {$P_{1}$};
\draw[trueedge] (n0.south) -- (n1.north);
\draw[trueedge] (n1.south) -- (n2.north);
\draw[falseedge] (n1.south) -- (n3.north);
\draw[trueedge] (n3.south) -- (n4.north);
\draw[falseedge] (n3.south) -- (n5.north);
\draw[trueedge] (n5.south) -- (n6.north);
\draw[falseedge] (n5.south) -- (n7.north);
\draw[falseedge] (n0.south) -- (n8.north);
\draw[trueedge] (n8.south) -- (n9.north);
\draw[trueedge] (n9.south) -- (n10.north);
\draw[trueedge] (n10.south) -- (n11.north);
\draw[trueedge] (n11.south) -- (n12.north);
\draw[falseedge] (n11.south) -- (n13.north);
\draw[trueedge] (n13.south) -- (n14.north);
\draw[trueedge] (n14.south) -- (n15.north);
\draw[trueedge] (n15.south) -- (n16.north);
\draw[trueedge] (n16.south) -- (n17.north);
\draw[falseedge] (n16.south) -- (n18.north);
\draw[trueedge] (n18.south) -- (n19.north);
\draw[trueedge] (n19.south) -- (n20.north);
\draw[trueedge] (n20.south) -- (n21.north);
\draw[trueedge] (n21.south) -- (n22.north);
\draw[trueedge] (n22.south) -- (n23.north);
\draw[falseedge] (n22.south) -- (n24.north);
\draw[falseedge] (n21.south) -- (n25.north);
\draw[trueedge] (n25.south) -- (n26.north);
\draw[trueedge] (n26.south) -- (n27.north);
\draw[falseedge] (n26.south) -- (n28.north);
\draw[falseedge] (n25.south) -- (n29.north);
\draw[falseedge] (n20.south) -- (n30.north);
\draw[falseedge] (n19.south) -- (n31.north);
\draw[trueedge] (n31.south) -- (n32.north);
\draw[falseedge] (n31.south) -- (n33.north);
\draw[falseedge] (n18.south) -- (n34.north);
\draw[falseedge] (n15.south) -- (n35.north);
\draw[falseedge] (n14.south) -- (n36.north);
\draw[falseedge] (n13.south) -- (n37.north);
\draw[falseedge] (n10.south) -- (n38.north);
\draw[falseedge] (n9.south) -- (n39.north);
\draw[trueedge] (n39.south) -- (n40.north);
\draw[trueedge] (n40.south) -- (n41.north);
\draw[trueedge] (n41.south) -- (n42.north);
\draw[trueedge] (n42.south) -- (n43.north);
\draw[trueedge] (n43.south) -- (n44.north);
\draw[trueedge] (n44.south) -- (n45.north);
\draw[trueedge] (n45.south) -- (n46.north);
\draw[trueedge] (n46.south) -- (n47.north);
\draw[falseedge] (n46.south) -- (n48.north);
\draw[falseedge] (n45.south) -- (n49.north);
\draw[falseedge] (n44.south) -- (n50.north);
\draw[falseedge] (n43.south) -- (n51.north);
\draw[trueedge] (n51.south) -- (n52.north);
\draw[trueedge] (n52.south) -- (n53.north);
\draw[trueedge] (n53.south) -- (n54.north);
\draw[falseedge] (n53.south) -- (n55.north);
\draw[trueedge] (n55.south) -- (n56.north);
\draw[trueedge] (n56.south) -- (n57.north);
\draw[falseedge] (n56.south) -- (n58.north);
\draw[trueedge] (n58.south) -- (n59.north);
\draw[falseedge] (n58.south) -- (n60.north);
\draw[falseedge] (n55.south) -- (n61.north);
\draw[trueedge] (n61.south) -- (n62.north);
\draw[falseedge] (n61.south) -- (n63.north);
\draw[trueedge] (n63.south) -- (n64.north);
\draw[falseedge] (n63.south) -- (n65.north);
\draw[falseedge] (n52.south) -- (n66.north);
\draw[falseedge] (n51.south) -- (n67.north);
\draw[trueedge] (n67.south) -- (n68.north);
\draw[trueedge] (n68.south) -- (n69.north);
\draw[falseedge] (n68.south) -- (n70.north);
\draw[trueedge] (n70.south) -- (n71.north);
\draw[falseedge] (n70.south) -- (n72.north);
\draw[trueedge] (n72.south) -- (n73.north);
\draw[falseedge] (n72.south) -- (n74.north);
\draw[falseedge] (n67.south) -- (n75.north);
\draw[falseedge] (n42.south) -- (n76.north);
\draw[trueedge] (n76.south) -- (n77.north);
\draw[trueedge] (n77.south) -- (n78.north);
\draw[trueedge] (n78.south) -- (n79.north);
\draw[trueedge] (n79.south) -- (n80.north);
\draw[falseedge] (n79.south) -- (n81.north);
\draw[falseedge] (n78.south) -- (n82.north);
\draw[falseedge] (n77.south) -- (n83.north);
\draw[falseedge] (n76.south) -- (n84.north);
\draw[trueedge] (n84.south) -- (n85.north);
\draw[trueedge] (n85.south) -- (n86.north);
\draw[trueedge] (n86.south) -- (n87.north);
\draw[trueedge] (n87.south) -- (n88.north);
\draw[falseedge] (n87.south) -- (n89.north);
\draw[trueedge] (n89.south) -- (n90.north);
\draw[falseedge] (n89.south) -- (n91.north);
\draw[trueedge] (n91.south) -- (n92.north);
\draw[falseedge] (n91.south) -- (n93.north);
\draw[falseedge] (n86.south) -- (n94.north);
\draw[trueedge] (n94.south) -- (n95.north);
\draw[falseedge] (n94.south) -- (n96.north);
\draw[trueedge] (n96.south) -- (n97.north);
\draw[trueedge] (n97.south) -- (n98.north);
\draw[falseedge] (n97.south) -- (n99.north);
\draw[falseedge] (n96.south) -- (n100.north);
\draw[falseedge] (n85.south) -- (n101.north);
\draw[trueedge] (n101.south) -- (n102.north);
\draw[falseedge] (n101.south) -- (n103.north);
\draw[trueedge] (n103.south) -- (n104.north);
\draw[falseedge] (n103.south) -- (n105.north);
\draw[falseedge] (n84.south) -- (n106.north);
\draw[trueedge] (n106.south) -- (n107.north);
\draw[trueedge] (n107.south) -- (n108.north);
\draw[trueedge] (n108.south) -- (n109.north);
\draw[falseedge] (n108.south) -- (n110.north);
\draw[falseedge] (n107.south) -- (n111.north);
\draw[falseedge] (n106.south) -- (n112.north);
\draw[trueedge] (n112.south) -- (n113.north);
\draw[falseedge] (n112.south) -- (n114.north);
\draw[falseedge] (n41.south) -- (n115.north);
\draw[trueedge] (n115.south) -- (n116.north);
\draw[trueedge] (n116.south) -- (n117.north);
\draw[falseedge] (n116.south) -- (n118.north);
\draw[trueedge] (n118.south) -- (n119.north);
\draw[falseedge] (n118.south) -- (n120.north);
\draw[trueedge] (n120.south) -- (n121.north);
\draw[trueedge] (n121.south) -- (n122.north);
\draw[trueedge] (n122.south) -- (n123.north);
\draw[trueedge] (n123.south) -- (n124.north);
\draw[falseedge] (n123.south) -- (n125.north);
\draw[falseedge] (n122.south) -- (n126.north);
\draw[falseedge] (n121.south) -- (n127.north);
\draw[falseedge] (n120.south) -- (n128.north);
\draw[falseedge] (n115.south) -- (n129.north);
\draw[trueedge] (n129.south) -- (n130.north);
\draw[trueedge] (n130.south) -- (n131.north);
\draw[falseedge] (n130.south) -- (n132.north);
\draw[trueedge] (n132.south) -- (n133.north);
\draw[trueedge] (n133.south) -- (n134.north);
\draw[trueedge] (n134.south) -- (n135.north);
\draw[falseedge] (n134.south) -- (n136.north);
\draw[falseedge] (n133.south) -- (n137.north);
\draw[falseedge] (n132.south) -- (n138.north);
\draw[falseedge] (n129.south) -- (n139.north);
\draw[trueedge] (n139.south) -- (n140.north);
\draw[trueedge] (n140.south) -- (n141.north);
\draw[trueedge] (n141.south) -- (n142.north);
\draw[trueedge] (n142.south) -- (n143.north);
\draw[trueedge] (n143.south) -- (n144.north);
\draw[falseedge] (n143.south) -- (n145.north);
\draw[falseedge] (n142.south) -- (n146.north);
\draw[falseedge] (n141.south) -- (n147.north);
\draw[falseedge] (n140.south) -- (n148.north);
\draw[falseedge] (n139.south) -- (n149.north);
\draw[trueedge] (n149.south) -- (n150.north);
\draw[trueedge] (n150.south) -- (n151.north);
\draw[trueedge] (n151.south) -- (n152.north);
\draw[falseedge] (n151.south) -- (n153.north);
\draw[falseedge] (n150.south) -- (n154.north);
\draw[falseedge] (n149.south) -- (n155.north);
\draw[falseedge] (n40.south) -- (n156.north);
\draw[trueedge] (n156.south) -- (n157.north);
\draw[trueedge] (n157.south) -- (n158.north);
\draw[trueedge] (n158.south) -- (n159.north);
\draw[falseedge] (n158.south) -- (n160.north);
\draw[trueedge] (n160.south) -- (n161.north);
\draw[trueedge] (n161.south) -- (n162.north);
\draw[trueedge] (n162.south) -- (n163.north);
\draw[trueedge] (n163.south) -- (n164.north);
\draw[trueedge] (n164.south) -- (n165.north);
\draw[trueedge] (n165.south) -- (n166.north);
\draw[falseedge] (n165.south) -- (n167.north);
\draw[falseedge] (n164.south) -- (n168.north);
\draw[falseedge] (n163.south) -- (n169.north);
\draw[trueedge] (n169.south) -- (n170.north);
\draw[trueedge] (n170.south) -- (n171.north);
\draw[falseedge] (n170.south) -- (n172.north);
\draw[falseedge] (n169.south) -- (n173.north);
\draw[falseedge] (n162.south) -- (n174.north);
\draw[falseedge] (n161.south) -- (n175.north);
\draw[falseedge] (n160.south) -- (n176.north);
\draw[falseedge] (n157.south) -- (n177.north);
\draw[falseedge] (n156.south) -- (n178.north);
\draw[falseedge] (n39.south) -- (n179.north);
\draw[falseedge] (n8.south) -- (n180.north);
\end{tikzpicture}
\end{minipage}\par\smallskip
{\footnotesize\textcolor{treeblue}{Solid: true}\quad\textcolor{treeorange}{Dashed: false}\quad $\mathbf B$: background\quad $\mathbf L$: lesion. NaN follows false.}
\end{minipage}}}
\captionsetup{type=figure}
\caption{\textbf{Both complete binary trees, with no collapsed branches.} The labels give each procedure's patient-grouped score; these separately selected all-patient refits have no additional held-out score. Repeated $P_j$ labels denote the same predicate at distinct positions. Every test and terminal decision is shown. The exact predicates are listed on page~2. Training weights and cutoff searches differ between procedures, so the accuracy gain is not attributed to tree size alone.
}
\label{fig:tree}
\end{minipage}}
\vskip 0.06in
]
\printAffiliationsAndNotice{}
\hypersetup{pdftitle={Compact Decision Programs over Frozen Neural Features for Lesion Segmentation},pdfsubject={Named-author research white paper}}
\section{Accuracy and Complete Decision Trees}
We pool each patient's scan counts, calculate $D_g=2\mathrm{TP}_g/(2\mathrm{TP}_g+\mathrm{FP}_g+\mathrm{FN}_g)$, and average over patients. The current procedure improves 28 patients, ties one, and lowers three versus the fixed mask: mean gain $0.031$, with descriptive 95\% patient-bootstrap interval $[0.010,0.052]$ (10,000 paired-patient resamples of saved predictions; unadjusted for repeated selection).

\begin{table}[h]
\caption{Cross-validated procedure scores and sizes of separately selected all-patient refits. The pictured trees have no additional held-out score. Storage excludes the CNN.}
\label{tab:performance}
\centering\small
\begin{tabular}{lrrr}
\toprule
Procedure & Dice & Tree nodes & Storage\\
\midrule
Fixed CNN mask & 0.734 & --- & ---\\
Earlier XGBoost & 0.758 & 21 & $\sim$250 B\\
Current XGBoost & \textbf{0.765} & 181 & $\sim$770 B\\
\bottomrule
\end{tabular}
\end{table}

The $0.007$ gain over the earlier procedure accompanies larger final decision logic (Fig.~\ref{fig:tree}). Fresh all-patient selection chooses depth~3, 16 rounds and FP32 cutoff~0.325. Its 32 source trees contain 480 nodes; conversion merges identical continuations into 63 shared nodes and 768 compact bytes in about half a second. The complete expanded tree has 90 tests and 91 terminal nodes. Construction checks pass; optional independent exhaustive qualification was not run. Native XGBoost~3.4.1 and C++23/CUDA complete evaluation and refitting in about 12 minutes after feature preparation.


\clearpage
\section{Patient-Grouped Learning}
We combine frozen LesionLocator features \citep{rokuss2025} with XGBoost \citep{chen2016}. RIDER Lung CT version~3 \citep{riderdata,zhao2024} supplies 32 patients: 30 repeat pairs and two single scans. Independent spreadsheet locations define an $80^3$ grid at 2-mm spacing. Two intended acquisitions lack usable location mapping and are withheld without contour-based correction. Missed foreground outside the region counts as false negatives; one patient's zero score is retained.

Each voxel has 54 inputs: 20 CT/location features, the pretrained mask, foreground probability, and 32 decoder activations. Probability uses the released five-model/eight-view ensemble; activations use fold zero's unmirrored pass. No neural weights are fitted locally. Unavailable continuous neural support is zero-filled.

Five outer folds hold out patients; five inner folds select the model and decision rule. Repeat scans always stay together. Eight weighted-cross-entropy candidates cross depths $\{1,2,3,5\}$ with $\{16,64\}$ rounds ($\eta=0.1$, $L_2=1$, 128 bins). Rules include argmax and 161 FP32 probability cutoffs. A second XGBoost stage offers four combiners, trained on patient-excluded predictions using three further folds.

Fitting retains all foreground and samples background at stride~32. The current procedure weights retained class-$c$ rows from patient $g$ by
\begin{equation}
 w_{gc}=\frac{n}{G}\frac{N_{gc}}{N_g n_{gc}},
\end{equation}
where $G,n$ count fitting patients and retained rows, and $N_g,N_{gc},n_{gc}$ are full patient, full class, and retained class counts. This preserves patient prevalence with equal patient mass. The earlier procedure uses global inverse-retention weights and 19 cutoffs. Within outer evaluation, selection and fitting weights exclude held-out patients.


\newpage
\section{Reading the Complete Trees}
Each solid branch means the test is true; each dashed branch means false. Terminal $B$ and $L$ denote background and lesion. Repeated $P_j$ labels identify the same test at different positions, not a collapsed branch.\par\smallskip
Here $x$ is normalized CT intensity; $u,s$ are local mean (2-mm radius) and standard deviation (4-mm radius); $v$ is intensity at the supplied point; $m,p$ are CNN mask and foreground probability; $a_k$ is decoder activation $k$.\par\smallskip
The constants below round-trip to stored FP32 thresholds. NaN follows the dashed branch. The node labels and line styles preserve meaning independently of color.\par
\smallskip{\fontsize{7}{8.5}\selectfont
\textbf{Current predicates}\par\smallskip
$P_{1}:\;m<1$\quad $P_{2}:\;p<0.0192318726$\par
$P_{3}:\;x<-0.537887037$\quad $P_{4}:\;v<0.147721395$\par
$P_{5}:\;x<-0.266500354$\quad $P_{6}:\;a_{28}<-0.0118003832$\par
$P_{7}:\;x<-0.419816613$\quad $P_{8}:\;a_{28}<-0.012050854$\par
$P_{9}:\;x<-1.07635784$\quad $P_{10}:\;u<-0.660179913$\par
$P_{11}:\;x<-1.4028399$\quad $P_{12}:\;a_{21}<-0.0135068875$\par
$P_{13}:\;s<1.26291883$\quad $P_{14}:\;a_{21}<-0.0102521647$\par
$P_{15}:\;v<0.207296506$\quad $P_{16}:\;a_{21}<-0.0117333885$\par
$P_{17}:\;s<1.18441689$\quad $P_{18}:\;u<-0.969308496$\par
$P_{19}:\;a_{22}<-0.0138405655$\quad $P_{20}:\;u<0.195440516$\par
$P_{21}:\;x<-0.76002878$\par
\textbf{Earlier thresholds}\par\smallskip
$\tau_{1}=-0.407344371$\quad $\tau_{2}=-0.711924314$\par
$\tau_{3}=-0.0117178345$\quad $\tau_{4}=-0.635912418$\par
$\tau_{5}=-0.0102527272$\quad $\tau_{6}=-0.785341203$\par
$\tau_{7}=-1.02144098$\par
}

\section{Scope and Availability}
The cohort informed earlier development; nested grouping does not make repeated use confirmatory \citep{cawley2010}. External checkpoint training overlap remains unverified. Independent patient evaluation is needed. The complete trees expose both procedures\textquotesingle{} decision logic while retaining the upstream CNN. Code, models and paper: \href{https://github.com/mingusb/rider-lesion-segmentation}{\texttt{github.com/mingusb/rider-lesion-segmentation}}.

\clearpage
\begin{thebibliography}{5}
\bibitem[Cawley \& Talbot(2010)]{cawley2010}
Cawley, G. C. and Talbot, N. L. C.
On over-fitting in model selection and subsequent selection bias in performance evaluation.
\emph{Journal of Machine Learning Research}, 11:2079--2107, 2010.
\url{https://www.jmlr.org/papers/v11/cawley10a.html}.

\bibitem[Chen \& Guestrin(2016)]{chen2016}
Chen, T. and Guestrin, C.
XGBoost: A scalable tree boosting system.
In \emph{Proceedings of the 22nd ACM SIGKDD International Conference on Knowledge Discovery and Data Mining}, pp.~785--794, 2016.
\doi{10.1145/2939672.2939785}.

\bibitem[Rokuss et~al.(2025)]{rokuss2025}
Rokuss, M., Kirchhoff, Y., Akbal, S., Kovacs, B., Roy, S., Ulrich, C., Wald, T., Rotkopf, L.~T., Schlemmer, H.-P., and Maier-Hein, K.
LesionLocator: Zero-shot universal tumor segmentation and tracking in 3D whole-body imaging.
In \emph{Proceedings of the IEEE/CVF Conference on Computer Vision and Pattern Recognition}, pp.~30872--30885, 2025.

\bibitem[Zhao et~al.(2015)]{riderdata}
Zhao, B., Schwartz, L.~H., Kris, M.~G., and Riely, G.~J.
Coffee-break lung CT collection with scan images reconstructed at multiple imaging parameters.
\emph{The Cancer Imaging Archive}, 2015. Dataset, version~3, updated June~25, 2024.
\doi{10.7937/K9/TCIA.2015.U1X8A5NR}.

\bibitem[Zhao et~al.(2024)]{zhao2024}
Zhao, B., Dercle, L., Yang, H., Riely, G.~J., Kris, M.~G., and Schwartz, L.~H.
Annotated test-retest dataset of lung cancer CT scan images reconstructed at multiple imaging parameters.
\emph{Scientific Data}, 11:1259, 2024.
\doi{10.1038/s41597-024-04085-3}.
\end{thebibliography}
\end{document}
